\oldpage[IV-3]{95}
\section{Preesquemas de Jacobson}
\label{section:IV.10}

Ya hemos tenido ocasión de observar \sref{IV.5.2.5} que incluso los
preesquemas \emph{excelentes} \sref{IV.7.8.5} no se comportan siempre como las
``variedades'' de la Geometría Algebraica clásica, especialmente en lo relativo
a cuestiones de dimensión; así, si $X$ es el espectro de un anillo de valoración
discreta completo, el conjunto de los puntos cerrados de $X$ (reducido a un solo
elemento) no es denso en todas partes en $X$, mientras que su complementario
(reducido a un solo elemento) es un abierto denso en todas partes de $X$ cuya
dimensión es cero, por tanto $<\dim(X)$. En esta sección examinamos condiciones
generales bajo las cuales no se producen tales fenómenos; como consecuencia,
las relaciones entre dimensión, codimensión, profundidad y coprofundidad en
estos preesquemas se comportan más satisfactoriamente en ciertos aspectos (10.6
y 10.8). Además, en los preesquemas considerados, el hecho de que el conjunto
de los puntos cerrados sea denso en todas partes (e incluso ``muy denso''
\sref{IV.10.1.3}) permite considerar únicamente estos puntos en muchas
demostraciones; recuperamos así el punto de vista clásico de las ``variedades
algebraicas'', que, desde nuestro punto de vista, son los conjuntos de puntos
cerrados de los \emph{preesquemas algebraicos sobre un cuerpo}, y establecemos
el vínculo entre el lenguaje de los esquemas y el de las ``variedades'' o
``espacios algebraicos'' de Serre (10.9 y 10.10).

\subsection{Subconjuntos muy densos de un espacio topológico}
\label{subsection:IV.10.1}

\begin{env}[10.1.1]
\label{IV.10.1.1}
Decimos que un subconjunto de un espacio topológico $X$ es \emph{cuasiconstruible} si es una unión finita de subconjuntos localmente cerrados de $X$.
Decimos que un subconjunto $T$ de $X$ es \emph{localmente cuasiconstruible} si, para todo $x\in X$, existe un entorno abierto $V$ de $x$ tal que $T\cap V$ es cuasiconstruible en $V$.
Está claro que todo subconjunto cuasiconstruible de $X$ es localmente cuasiconstruible; el recíproco es cierto si $X$ es cuasicompacto.
El razonamiento de \sref[$0_{\mathrm{III}}$]{0.9.1.3}, en el que se suprime la palabra ``retrocompacto'', muestra que el conjunto de los subconjuntos cuasiconstruibles (resp.\ localmente cuasiconstruibles) de $X$, que denotaremos por $\mathfrak{Qc}(X)$ (resp.\ $\mathfrak{Lqc}(X)$), es estable por unión finita, intersección finita y paso al complementario.
Si $f:X\to Y$ es una aplicación continua, se deduce inmediatamente de estas definiciones que, para todo subconjunto cuasiconstruible (resp.\ localmente cuasiconstruible) $Z$ de $Y$, $f^{-1}(Z)$ es cuasiconstruible (resp.\ localmente cuasiconstruible) en $X$.

Los subconjuntos construibles (resp.\ localmente construibles) de $X$ son evidentemente cuasiconstruibles (resp.\ localmente cuasiconstruibles); el recíproco es cierto cuando $X$ es noetheriano (resp.\ localmente noetheriano).

En lo que sigue, denotaremos respectivamente por $\mathfrak{O}(X)$, $\mathfrak{F}(X)$, $\mathfrak{Lf}(X)$, $\mathfrak{C}(X)$, $\mathfrak{Lc}(X)$ el conjunto de los subconjuntos de $X$ que son respectivamente abiertos, cerrados, localmente cerrados, construibles y localmente construibles.
\end{env}

\oldpage[IV-3]{96}

\begin{proposition}[10.1.2]
\label{IV.10.1.2}
Sea $X$ un espacio topológico y $X_0$ un subconjunto de $X$.
Las condiciones siguientes son equivalentes:
\begin{enumerate}[label=\emph{\alph*)}]
  \item Para todo subconjunto localmente cerrado $Z\neq\varnothing$ de $X$, se tiene $Z\cap X_0\neq\varnothing$.
  \item[\rm a$'$)] Para todo subconjunto cerrado $Z$ de $X$, se tiene $Z=\overline{Z\cap X_0}$.
  \item Para todo subconjunto localmente cuasiconstruible $Z\neq\varnothing$ de $X$, se tiene $Z\cap X_0\neq\varnothing$.
  \item[\rm b$'$)] Para todo subconjunto localmente cuasiconstruible $Z$ de $X$, se tiene $Z\subset\overline{Z\cap X_0}$ (dicho de otro modo, $Z\cap X_0$ es denso en $Z$).
  \item La aplicación $U\mapsto U\cap X_0$ de $\mathfrak{O}(X)$ en $\mathfrak{O}(X_0)$ es inyectiva (y por tanto biyectiva).
  \item[\rm c$'$)] La aplicación $F\mapsto F\cap X_0$ de $\mathfrak{F}(X)$ en $\mathfrak{F}(X_0)$ es inyectiva (y por tanto biyectiva).
  \item[\rm c$''$)] La aplicación $Z\mapsto Z\cap X_0$ de $\mathfrak{Qc}(X)$ en $\mathfrak{Qc}(X_0)$ es inyectiva (y por tanto biyectiva).
  \item[\rm c$'''$)] La aplicación $Z\mapsto Z\cap X_0$ de $\mathfrak{Lqc}(X)$ en $\mathfrak{Lqc}(X_0)$ es inyectiva.
\end{enumerate}
Además, cuando se satisfacen estas condiciones, la aplicación $Z\mapsto Z\cap X_0$ de $\mathfrak{Lqc}(X)$ en $\mathfrak{Lqc}(X_0)$ es biyectiva.
\end{proposition}

\begin{proof}
Observemos que las afirmaciones de sobreyectividad de $c)$ y $c')$ son triviales; significan que todo subconjunto localmente cerrado de $X_0$ es la traza sobre $X_0$ de un subconjunto localmente cerrado de $X$, y que la aplicación $\mathfrak{Qc}(X)\to\mathfrak{Qc}(X_0)$ definida en $c'')$ también es sobreyectiva.

Demostraremos las implicaciones
\[
  c''')\Rightarrow c'')\Rightarrow c')\Rightarrow c)\Rightarrow b)\Rightarrow b')\Rightarrow a')\Rightarrow a)\Rightarrow c''').
\]

Las tres primeras son triviales.
Para ver que $c)$ implica $b)$, observemos primero que $c)$ significa que $X_0$ es \emph{denso} en $X$.
Sustituyendo respectivamente $X$ y $X_0$ por $U$ y $U\cap X_0$, podemos por tanto limitarnos al caso en que $Z$ es localmente cerrado en $X$, y entonces $Z=V\cap\complement W$, donde $V$ y $W$ son abiertos de $X$; la hipótesis $Z\neq\varnothing$ significa que $V\not\subset W$, o equivalentemente $V\cup W\neq W$.
En virtud de $c)$, se tiene $(V\cup W)\cap X_0\neq W\cap X_0$, luego $V\cap X_0\not\subset W$, y por consiguiente $(V\cap\complement W)\cap X_0\neq\varnothing$.

Para ver que $b)$ implica $b')$, basta aplicar $b)$ a $Z\cap U$, donde $U$ es un entorno abierto de un punto arbitrario de $Z$.
Como $Z\cap\overline{X_0}\subset\overline{Z\cap X_0}$, es trivial que $b')$ implica $a')$.
Para mostrar que $a')$ implica $a)$, observemos que, si $Z$ es localmente cerrado en $X$, podemos escribir $Z=F-F'$, donde $F$ y $F'$ son cerrados de $X$ y $F'\subset F$; de donde $Z\cap X_0=(F\cap X_0)-(F'\cap X_0)$.
Si $Z\cap X_0=\varnothing$, entonces, en virtud de $a')$, tendríamos $F\cap X_0=F'\cap X_0$, luego $F=F'$, es decir, $Z=\varnothing$.

Para ver que $a)$ implica $c''')$, basta mostrar que, si $Z\neq\varnothing$ es localmente cuasiconstruible, entonces $Z\cap X_0\neq\varnothing$: en efecto, la relación $Z'\cap X_0=Z''\cap X_0$ equivale a $((Z'\cup Z'')-(Z'\cap Z''))\cap X_0=\varnothing$.
Dicho de otro modo, basta probar que $a)$ implica $b)$; en efecto, sustituyendo respectivamente $X$ y $X_0$ por $U$ y $U\cap X_0$, nos reducimos al caso en que $Z$ es localmente cerrado en $X$, de donde se sigue la conclusión.

Queda por mostrar que la aplicación $\mathfrak{Lqc}(X)\to\mathfrak{Lqc}(X_0)$ es sobreyectiva.
Sea pues $Z_0$ un subconjunto localmente cuasiconstruible de $X_0$: existe un recubrimiento $(V_\alpha)$ de $X_0$ por abiertos de $X_0$ tal que $Z_0\cap V_\alpha$ es cuasiconstruible en $V_\alpha$ (y
\oldpage[IV-3]{97}
por tanto también en $X_0$).
En virtud de $c)$, para todo $\alpha$ existe un único abierto $U_\alpha$ de $X$ tal que $X_0\cap U_\alpha=V_\alpha$, y en virtud de $c'')$ un único subconjunto cuasiconstruible $Z_\alpha$ de $X$ tal que $Z_0\cap V_\alpha=X_0\cap Z_\alpha$.
Si $\alpha$ y $\beta$ son dos índices cualesquiera, se tiene $Z_\beta\cap U_\alpha\cap X_0=U_\alpha\cap Z_0\cap V_\beta=V_\alpha\cap Z_0\cap V_\beta=Z_\alpha\cap X_0\cap V_\beta\subset Z_\alpha\cap X_0$; como $Z_\beta\cap U_\alpha$ y $Z_\alpha$ son cuasiconstruibles en $X$, se deduce de $c'')$ que $Z_\beta\cap U_\alpha\subset Z_\alpha$.
Si ponemos $Z=\bigcup_\alpha Z_\alpha$, entonces $Z\cap U_\alpha=Z_\alpha$ para todo $\alpha$; además, puesto que los $V_\alpha$ recubren $X_0$, se deduce de $c)$ que los $U_\alpha$ recubren $X$, y vemos por tanto que $Z$ es localmente cuasiconstruible en $X$ y que $Z_0=Z\cap X_0$.
\end{proof}

\begin{definition}[10.1.3]
\label{IV.10.1.3}
Cuando un subconjunto $X_0$ de un espacio topológico satisface las condiciones equivalentes de \sref{IV.10.1.2}, decimos que $X_0$ es \emph{muy denso} en $X$.
\end{definition}

Ya hemos visto en el curso de la demostración de \sref{IV.10.1.2} que entonces $X_0$ es denso en $X$.

\begin{corollary}[10.1.4]
\label{IV.10.1.4}
Si $X_0$ es muy denso en $X$ y $U$ es un abierto de $X$, $U\cap X_0$ es muy denso en $U$.
Recíprocamente, si $(U_\alpha)$ es un recubrimiento abierto de $X$ tal que $U_\alpha\cap X_0$ es muy denso en $U_\alpha$ para todo $\alpha$, entonces $X_0$ es muy denso en $X$.
\end{corollary}

Como todo subconjunto localmente cerrado de $U$ es localmente cerrado en $X$, la primera afirmación se deduce del criterio $a)$ de \sref{IV.10.1.2}; y lo mismo vale para la segunda, pues si $Z\neq\varnothing$ es localmente cerrado en $X$, $Z\cap U_\alpha$ es localmente cerrado en $U_\alpha$ para todo $\alpha$, y $Z\cap U_\alpha\neq\varnothing$ para al menos un $\alpha$.

\subsection{Cuasihomeomorfismos}
\label{subsection:IV.10.2}

\begin{proposition}[10.2.1]
\label{IV.10.2.1}
Sean $X_0$, $X$ dos espacios topológicos y $f:X_0\to X$ una aplicación continua.
Las condiciones siguientes son equivalentes:
\begin{enumerate}[label=\emph{\alph*)}]
  \item La aplicación $U\mapsto f^{-1}(U)$ de $\mathfrak{O}(X)$ en $\mathfrak{O}(X_0)$ es biyectiva.
  \item[\rm a$'$)] La aplicación $F\mapsto f^{-1}(F)$ de $\mathfrak{F}(X)$ en $\mathfrak{F}(X_0)$ es biyectiva.
  \item La topología de $X_0$ es la imagen inversa por $f$ de la de $X$, y $f(X_0)$ es muy denso \sref{IV.10.1.3} en $X$.
  \item El funtor $\sh{F}\mapsto f^*(\sh{F})$ de la categoría de haces de conjuntos (resp.\ haces de grupos abelianos) sobre $X$ en la categoría de haces de conjuntos (resp.\ haces de grupos abelianos) sobre $X_0$ es una equivalencia de categorías.
\end{enumerate}
\end{proposition}

\begin{proof}
Está claro que $a)$ y $a')$ son equivalentes, y que $a)$ implica que la topología de $X_0$ es la imagen inversa por $f$ de la de $X$.
Por otra parte, si $f(X_0)$ no es muy denso en $X$, existen dos abiertos distintos $U_1$, $U_2$ de $X$ tales que $U_1\cap f(X_0)=U_2\cap f(X_0)$, y por tanto $f^{-1}(U_1)=f^{-1}(U_2)$, lo que termina de mostrar que $a)$ implica $b)$.
Recíprocamente, la condición $b)$ implica que las aplicaciones $U\mapsto U\cap f(X_0)$ de $\mathfrak{O}(X)$ en $\mathfrak{O}(f(X_0))$ y $V\mapsto f^{-1}(V)$ de $\mathfrak{O}(f(X_0))$ en $\mathfrak{O}(X_0)$ son biyectivas, y por tanto también lo es su compuesta $U\mapsto f^{-1}(U)$.

Para ver que $a)$ implica $c)$, basta aplicar la definición de $f^*(\sh{F})$ \sref[0]{0.3.7.1} y los axiomas de los haces: $a)$ significa que, para todo abierto $U$ de $X$, la aplicación canónica $\Gamma(U,\sh{F})\to\Gamma(f^{-1}(U),f^*(\sh{F}))$ es una biyección funtorial en $\sh{F}$, de donde $c)$.
Queda por mostrar que $c)$ implica $b)$.

\oldpage[IV-3]{98}
Supongamos primero que $f(X_0)$ no es muy denso en $X$; existen entonces dos cerrados distintos $Y\supset Y'$ de $X$ tales que $Y\cap f(X_0)=Y'\cap f(X_0)$.
Sea $\sh{F}$ (resp.\ $\sh{F}'$) el haz de grupos abelianos sobre $X$ que es la imagen directa por la inyección canónica $Y\to X$ (resp.\ $Y'\to X$) del haz simple asociado al prehaz constante $\mathbf{Z}$ sobre $Y$ (resp.\ $Y'$); la definición de $f^*$ \sref[0]{0.3.7.1} muestra que $f^*(\sh{F})$ es isomorfo a $f^*(\sh{F}')$, pero $\sh{F}$ no es isomorfo a $\sh{F}'$, y por tanto no se satisface la condición $c)$.
(Obsérvese que el funtor $f^*$ ni siquiera es entonces \emph{fiel}, pues si $u$ y $v$ son el automorfismo identidad y el endomorfismo nulo de $\sh{F}/\sh{F}'$, $f^*(u)$ y $f^*(v)$ son iguales.)

Mostremos ahora que $c)$ implica que la topología de $X_0$ es la imagen inversa por $f$ de la de $X$.
Observemos que, si la condición $c)$ se satisface para la categoría de haces de conjuntos, también se satisface para la categoría de haces de grupos abelianos, pues se deduce inmediatamente de las definiciones \sref[$0_{\mathrm{III}}$]{0.8.2.5} que la categoría de los \emph{objetos en grupos abelianos} de la categoría de haces de conjuntos no es otra que la categoría de haces de grupos abelianos.
Bastará pues probar nuestra afirmación suponiendo satisfecha $c)$ para las categorías de haces de grupos abelianos sobre $X$ y $X_0$.
Ahora bien, si $\mathbf{Z}_X$ denota el haz simple sobre $X$ asociado al prehaz constante $\mathbf{Z}$, existe un isomorfismo canónico $\Gamma(X,\sh{F})\xrightarrow{\sim}\operatorname{Hom}(\mathbf{Z}_X,\sh{F})$ funtorial en $\sh{F}$ (por el mismo razonamiento que en \sref[0]{0.5.1.1}); como por definición \sref[0]{0.3.7.1} está claro que $f^*(\mathbf{Z}_X)=\mathbf{Z}_{X_0}$, la hipótesis de que $f^*$ es una equivalencia significa que el homomorfismo canónico $\Gamma(X,\sh{F})\to\Gamma(X_0,f^*(\sh{F}))$ es \emph{biyectivo} para todo haz de grupos abelianos $\sh{F}$ sobre $X$.
Sea ahora $Y_0$ un cerrado de $X_0$; sea $\sh{F}_0$ el haz sobre $X_0$ que es la imagen directa por la inyección canónica $Y_0\to X_0$ del haz simple asociado al prehaz constante $\mathbf{Z}$ sobre $Y_0$.
Como $f^*$ es una equivalencia, $\sh{F}_0$ es isomorfo a un haz de la forma $f^*(\sh{F})$, donde $\sh{F}$ es un haz de grupos abelianos sobre $X$.
Consideremos entonces una sección $s_0$ de $\sh{F}_0$ sobre $X_0$ tal que $(s_0)_{x_0}=1_{x_0}$ si $x_0\in Y_0$, y $(s_0)_{x_0}=0_{x_0}$ si $x_0\notin Y_0$.
Las observaciones precedentes muestran que existe una sección $s$ de $\sh{F}$ sobre $X$ tal que $(s_0)_{x_0}=s_{f(x_0)}$ para todo $x_0\in X_0$ (estando identificadas canónicamente las fibras $(\sh{F}_0)_{x_0}$ y $\sh{F}_{f(x_0)}$ \sref[0]{0.3.7.1}); esto significa que $Y_0$ es la imagen inversa por $f$ del conjunto $Y$ de los $x\in X$ tales que $s_x\neq 0_x$, y $Y$ es cerrado en $X$.
\end{proof}

\begin{definition}[10.2.2]
\label{IV.10.2.2}
Cuando una aplicación continua $f:X_0\to X$ satisface las condiciones equivalentes de \sref{IV.10.2.1}, decimos que $f$ es un \emph{cuasihomeomorfismo} de $X_0$ en $X$.
\end{definition}

En virtud de \sref{IV.10.2.1},~$b)$, decir que un subconjunto $X_0$ de un espacio topológico $X$ es muy denso en $X$ significa que la inyección canónica $X_0\to X$ es un cuasihomeomorfismo.

\begin{corollary}[10.2.3]
\label{IV.10.2.3}
La compuesta de dos cuasihomeomorfismos es un cuasihomeomorfismo.
\end{corollary}

Esto se deduce inmediatamente de \sref{IV.10.2.1},~$a)$.

\begin{corollary}[10.2.4]
\label{IV.10.2.4}
Sea $f:X\to Y$ un cuasihomeomorfismo, $Y'$ un subconjunto localmente cuasiconstruible de $Y$ y $X'=f^{-1}(Y')$; entonces la restricción $f'=f|X':X'\to Y'$ es un cuasihomeomorfismo.
\end{corollary}

\begin{proof}
Está claro, en virtud de \sref{IV.10.2.1},~$b)$, que la topología inducida sobre $X'$ por la de $X$ es la imagen inversa por $f'$ de la topología inducida sobre $Y'$ por la de $Y$.
Por otra parte,
\oldpage[IV-3]{99}
sea $Z\neq\varnothing$ un subconjunto cerrado \emph{en} $Y'$, y sea $z\in Z$; existe un entorno abierto $U$ de $z$ en $Y$ tal que $U\cap Y'$ es la unión de un número finito de cerrados $Y'_j$ de $U$; si $j$ es un índice tal que $z\in Y'_j$, $Z\cap Y'_j$ es por tanto cerrado en $U$.
Como $f(X)\cap U$ es muy denso en $U$ \sref{IV.10.1.4}, $Z\cap Y'_j\cap f(X)$ no es vacío \sref{IV.10.1.2},~$a)$, ni, \emph{a fortiori}, $Z\cap f(X)$; pero, como $Z\subset Y'$, se tiene $Z\cap f(X)=Z\cap f'(X')$; el criterio \sref{IV.10.1.2},~$a)$ muestra que $f'(X')$ es muy denso en $Y'$, y concluimos por \sref{IV.10.2.1},~$b)$.
\end{proof}

\begin{corollary}[10.2.5]
\label{IV.10.2.5}
Sea $f:X\to Y$ una aplicación continua y $(V_\alpha)$ un recubrimiento abierto de $Y$.
Si, para todo $\alpha$, la restricción $f^{-1}(V_\alpha)\to V_\alpha$ de $f$ es un cuasihomeomorfismo, entonces $f$ es un cuasihomeomorfismo.
\end{corollary}

Esto se deduce inmediatamente del criterio \sref{IV.10.2.1},~$b)$ y de \sref{IV.10.1.4}.

\begin{corollary}[10.2.6]
\label{IV.10.2.6}
Sea $f:X\to Y$ un cuasihomeomorfismo, $Y'$ un subconjunto localmente cuasiconstruible de $Y$ y $X'=f^{-1}(Y')$.
Para que $Y'$ sea cuasicompacto (resp.\ noetheriano, resp.\ retrocompacto en $Y$), es necesario y suficiente que $X'$ sea cuasicompacto (resp.\ noetheriano, resp.\ retrocompacto en $X$).
\end{corollary}

\begin{proof}
Probemos las dos primeras afirmaciones; en virtud de \sref{IV.10.2.4}, podemos limitarnos al caso en que $Y'=Y$.
Decir que $X$ es cuasicompacto (resp.\ noetheriano) significa que, para toda familia filtrante creciente $(U_\alpha)$ de $\mathfrak{O}(X)$ que tenga a $X$ como elemento máximo (resp.\ para toda familia filtrante creciente $(U_\alpha)$ de $\mathfrak{O}(X)$), existe $\gamma$ tal que $U_\alpha=U_\gamma$ para $\alpha\geq\gamma$.
Como $U\mapsto f^{-1}(U)$ es una biyección de $\mathfrak{O}(Y)$ sobre $\mathfrak{O}(X)$, nuestra afirmación se deduce inmediatamente de la observación precedente.

Los abiertos cuasicompactos de $X$ son por tanto los conjuntos de la forma $f^{-1}(U)$, donde $U$ es un abierto cuasicompacto de $Y$, en virtud de \sref{IV.10.2.1},~$a)$ y de lo que precede.
Para que $X'$ sea retrocompacto en $X$, es necesario y suficiente que, para todo abierto cuasicompacto $U$ de $Y$, $f^{-1}(U)\cap X'=f^{-1}(U\cap Y')$ sea cuasicompacto \sref[$0_{\mathrm{III}}$]{0.9.1.1}; la primera parte de la demostración muestra que esto equivale a decir que $U\cap Y'$ es cuasicompacto para todo abierto cuasicompacto $U$, es decir, que $Y'$ es retrocompacto en $Y$.
\end{proof}

\begin{proposition}[10.2.7]
\label{IV.10.2.7}
Sea $f:X\to Y$ un cuasihomeomorfismo.
Entonces la aplicación $Z\mapsto f^{-1}(Z)$ de $\mathfrak{P}(Y)$ en $\mathfrak{P}(X)$ define por restricción las biyecciones (cf.\ \sref{IV.10.1.1} para la notación):
\[
  \mathfrak{O}(Y)\xrightarrow{\sim}\mathfrak{O}(X)
\]
\[
  \mathfrak{F}(Y)\xrightarrow{\sim}\mathfrak{F}(X)
\]
\[
  \mathfrak{Lf}(Y)\xrightarrow{\sim}\mathfrak{Lf}(X)
\]
\[
  \mathfrak{Qc}(Y)\xrightarrow{\sim}\mathfrak{Qc}(X)
\]
\[
  \mathfrak{Lqc}(Y)\xrightarrow{\sim}\mathfrak{Lqc}(X)
\]
\[
  \mathfrak{C}(Y)\xrightarrow{\sim}\mathfrak{C}(X)
\]
\[
  \mathfrak{Lc}(Y)\xrightarrow{\sim}\mathfrak{Lc}(X).
\]
\end{proposition}

\begin{proof}
\oldpage[IV-3]{100}
Para las dos primeras, esto no es sino la definición \sref{IV.10.2.2}; como la topología de $X$ es la imagen inversa por $f$ de la de $f(X)$, las cinco aplicaciones restantes, en las que se sustituye $Y$ por $f(X)$, son biyectivas.
Basta pues (por \sref{IV.10.2.1},~$b)$) limitarnos al caso en que $X$ es un subespacio muy denso de $Y$, y ya se ha probado \sref{IV.10.1.2} que las aplicaciones $\mathfrak{Lf}(Y)\to\mathfrak{Lf}(X)$, $\mathfrak{Qc}(Y)\to\mathfrak{Qc}(X)$, $\mathfrak{Lqc}(Y)\to\mathfrak{Lqc}(X)$ son biyectivas.
Como todo subconjunto localmente construible es localmente cuasiconstruible, las aplicaciones $\mathfrak{C}(Y)\to\mathfrak{C}(X)$ y $\mathfrak{Lc}(Y)\to\mathfrak{Lc}(X)$ son inyectivas; además, para todo abierto $U\subset Y$, la restricción $f^{-1}(U)\to U$ de $f$ es un cuasihomeomorfismo \sref{IV.10.2.4}, por lo que, si mostramos que $\mathfrak{C}(Y)\to\mathfrak{C}(X)$ es sobreyectiva, lo mismo será cierto para $\mathfrak{Lc}(Y)\to\mathfrak{Lc}(X)$.
Pero, en virtud de \sref{IV.10.2.6}, todo abierto retrocompacto $Z$ de $X$ es de la forma $f^{-1}(T)$, donde $T$ es un abierto retrocompacto de $Y$; esto prueba evidentemente la sobreyectividad de $\mathfrak{C}(Y)\to\mathfrak{C}(X)$.
\end{proof}

\begin{remarks}[10.2.8]
\label{IV.10.2.8}
(i) Vimos en la demostración de \sref{IV.10.2.1} que, si $f$ es un cuasihomeomorfismo, la aplicación canónica
\begin{equation}
\label{IV.10.2.8.1}
\tag{10.2.8.1}
\Gamma(Y,\sh{F})\to\Gamma(X,f^*(\sh{F}))
\end{equation}
es un isomorfismo funtorial en $\sh{F}$ en la categoría de haces de grupos abelianos sobre $Y$.
Como $f^*$ es exacto en esta categoría \sref[0]{0.3.7.2}, esto implica que los homomorfismos canónicos $(T,3.2.2)$
\[
  \mathrm{H}^i(Y,\sh{F})\xrightarrow{\sim}\mathrm{H}^i(X,f^*(\sh{F}))
\]
son \emph{biyectivos} para todo $i$.

(ii) Si $f:X\to Y$ es una aplicación continua y $\sh{F}$, $\sh{G}$ son dos haces de grupos abelianos sobre $Y$, se tiene $f^*(\sh{F}\otimes_{\mathbf{Z}_Y}\sh{G})=f^*(\sh{F})\otimes_{\mathbf{Z}_X}f^*(\sh{G})$, salvo isomorfismo canónico \sref[0]{0.4.3.3}.
Concluimos que, si $f$ es un cuasihomeomorfismo, $f^*$ también es una \emph{equivalencia} entre la categoría de haces de anillos sobre $Y$ y la categoría de haces de anillos sobre $X$; el dato de una estructura de espacio anillado sobre $Y$ equivale por tanto al de una estructura de espacio anillado sobre $X$.

Dados dos espacios anillados $(X,\sh{O}_X)$, $(Y,\sh{O}_Y)$, diremos que un morfismo de espacios anillados $f=(\psi,\theta):(X,\sh{O}_X)\to(Y,\sh{O}_Y)$ es un \emph{cuasiisomorfismo} si $\psi$ es un cuasihomeomorfismo de $X$ en $Y$ y $\theta^\sharp:\psi^*(\sh{O}_Y)\to\sh{O}_X$ es un isomorfismo de haces de anillos; cuando esto ocurre, el espacio anillado $(X,\sh{O}_X)$ queda enteramente determinado, salvo isomorfismo, por $(Y,\sh{O}_Y)$, el espacio $X$ y el cuasihomeomorfismo $\psi$ (que puede elegirse arbitrariamente).
Cuando $f$ es un cuasiisomorfismo de espacios anillados, el funtor
\[
  \sh{F}\mapsto f^*(\sh{F})
\]
es una \emph{equivalencia} entre la categoría de $\sh{O}_Y$-Módulos y la de $\sh{O}_X$-Módulos, pues aquí $f^*(\sh{F})$ se identifica canónicamente con $\psi^*(\sh{F})$.
Obtenemos, por ejemplo, isomorfismos de bi-$\partial$-funtores
\[
  \operatorname{Ext}^i_{\sh{O}_Y}(\sh{F},\sh{G})\xrightarrow{\sim}\operatorname{Ext}^i_{\sh{O}_X}(f^*(\sh{F}),f^*(\sh{G})).
\]

En general, puede decirse que las construcciones usuales de la teoría de haces y del álgebra homológica, efectuadas sobre el espacio anillado $Y$ o sobre el espacio anillado $X$, son equivalentes.
\end{remarks}

\subsection{Espacios de Jacobson}
\label{subsection:IV.10.3}

\oldpage[IV-3]{101}

\begin{definition}[10.3.1]
\label{IV.10.3.1}
Decimos que un espacio topológico $X$ es un \emph{espacio de Jacobson} si el conjunto $X_0$ de los puntos cerrados de $X$ es muy denso en $X$ (dicho de otro modo, si la inyección canónica $X_0\to X$ es un cuasihomeomorfismo).
\end{definition}

Esto significa por tanto \sref{IV.10.1.2} que todo subconjunto localmente cerrado (o solamente localmente cuasiconstruible) $Z\neq\varnothing$ de $X$ contiene un punto cerrado de $X$, o equivalentemente que \emph{todo subconjunto cerrado de $X$ es la adherencia del conjunto de sus puntos cerrados}.

\begin{proposition}[10.3.2]
\label{IV.10.3.2}
Sea $X$ un espacio de Jacobson y $Z$ un subconjunto localmente cuasiconstruible de $X$; entonces el subespacio $Z$ es un espacio de Jacobson, y para que un punto de $Z$ sea cerrado en $Z$ es necesario y suficiente que sea cerrado en $X$.
\end{proposition}

\begin{proof}
Si $X_0$ es el conjunto de los puntos cerrados de $X$, $Z\cap X_0$ es muy denso en $Z$ en virtud de \sref{IV.10.2.4}, aplicado a la inyección $i:X_0\to X$; como el conjunto $Z_0$ de los puntos cerrados de $Z$ contiene evidentemente a $Z\cap X_0$, es \emph{a fortiori} muy denso en $Z$, luego $Z$ es un espacio de Jacobson.
Mostremos ahora que se tiene $Z_0=Z\cap X_0$; sea $x\in Z$ un punto cerrado en $Z$; sea $\overline{\{x\}}$ su adherencia en $X$. Entonces $\overline{\{x\}}\cap Z=\{x\}$ es un subconjunto localmente cuasiconstruible de $X$, y como su intersección con $X_0$ no es vacía \sref{IV.10.1.2}, se tiene $x\in X_0$.
\end{proof}

\begin{proposition}[10.3.3]
\label{IV.10.3.3}
Sea $X$ un espacio topológico y $(U_\alpha)$ un recubrimiento abierto de $X$.
Para que $X$ sea un espacio de Jacobson, es necesario y suficiente que cada uno de los subespacios $U_\alpha$ sea un espacio de Jacobson.
\end{proposition}

\begin{proof}
La condición es necesaria en virtud de \sref{IV.10.3.2}.
Recíprocamente, mostremos primero que la hipótesis de que los $U_\alpha$ son espacios de Jacobson implica que, para que un punto $x\in U_\alpha$ sea cerrado en $X$, basta que sea cerrado en $U_\alpha$.
Basta en efecto ver que esta condición implica que $x$ también es cerrado en cada uno de los $U_\beta$ que lo contienen; pero $U_\alpha\cap U_\beta$ es abierto en $U_\alpha$, luego $x$ es cerrado en $U_\alpha\cap U_\beta$, y en virtud de \sref{IV.10.3.2}, $x$ también es cerrado en $U_\beta$, lo que termina la demostración.
\end{proof}

\subsection{Preesquemas de Jacobson y anillos de Jacobson}
\label{subsection:IV.10.4}

\begin{definition}[10.4.1]
\label{IV.10.4.1}
Decimos que un preesquema $X$ es un \emph{preesquema de Jacobson} si el espacio topológico subyacente a $X$ es un espacio de Jacobson.
Decimos que un anillo $A$ es un \emph{anillo de Jacobson} si $\operatorname{Spec}(A)$ es un preesquema de Jacobson.
\end{definition}

Todo cerrado de $X=\operatorname{Spec}(A)$ es de la forma $Z=V(\mathfrak{a})$, donde $\mathfrak{a}$ es un ideal igual a su radical, y el conjunto $X_0$ de los puntos cerrados de $X$ es el conjunto de los ideales maximales de $A$; decir que $Z\cap X_0$ es denso en $Z$ significa por tanto que $\mathfrak{a}$ es una \emph{intersección de ideales maximales} \sref[I]{I.1.1.4}; como $\mathfrak{a}$ es una intersección de ideales primos, esto equivale a decir que todo ideal primo de $A$ es una intersección de ideales maximales; en virtud de \sref{IV.10.3.1} y \sref{IV.10.1.2}, la definición usual de los anillos de Jacobson (Bourbaki, \emph{Alg.\ comm.}, cap.~V, \S\,3, n$^\circ$\,4, def.~1) coincide así con la definición \sref{IV.10.4.1}.

\begin{proposition}[10.4.2]
\label{IV.10.4.2}
\oldpage[IV-3]{102}
Sea $X$ un preesquema y $(U_\alpha)$ un recubrimiento de $X$ formado por abiertos afines.
Para que $X$ sea un preesquema de Jacobson, es necesario y suficiente que los anillos $\Gamma(U_\alpha,\sh{O}_X)$ de los $U_\alpha$ sean anillos de Jacobson.
\end{proposition}

Esto se deduce de \sref{IV.10.3.3} y de la definición \sref{IV.10.4.1}.

\begin{env}[10.4.3]
\label{IV.10.4.3}
Un preesquema discreto es un preesquema de Jacobson; un anillo artiniano es por tanto un anillo de Jacobson.
Todo anillo de ideales principales que tenga una infinidad de ideales maximales (por ejemplo $\mathbf{Z}$) es un anillo de Jacobson; un anillo local noetheriano $A$ es un anillo de Jacobson si y solo si su ideal maximal es su único ideal primo, es decir, si $A$ es artiniano.
Todo subpreesquema de un preesquema de Jacobson es un preesquema de Jacobson en virtud de \sref{IV.10.3.2}.
\end{env}

\begin{proposition}[10.4.4]
\label{IV.10.4.4}
Sea $B$ un anillo íntegro.
Las condiciones siguientes son equivalentes:
\begin{enumerate}[label=\emph{\alph*)}]
  \item Existe $f\neq 0$ en $B$ tal que $B_f$ es un cuerpo.
  \item El cuerpo de fracciones de $B$ es una $B$-álgebra de tipo finito.
  \item Existe un cuerpo $K$ que contiene a $B$ y que es una $B$-álgebra de tipo finito.
  \item El punto genérico de $\operatorname{Spec}(B)$ es aislado.
\end{enumerate}
\end{proposition}

\begin{proof}
Está claro que $d)$ equivale a $a)$, pues $d)$ significa que existe $f\neq 0$ en $B$ tal que $D(f)$ se reduce al punto genérico de $\operatorname{Spec}(B)$.
Es trivial que $a)$ implica $b)$ y que $b)$ implica $c)$.
Finalmente, $c)$ implica $a)$, en virtud de Bourbaki, \emph{Alg.\ comm.}, cap.~V, \S\,3, n$^\circ$\,1, cor.~2 del th.~1.
\end{proof}

\begin{proposition}[10.4.5]
\label{IV.10.4.5}
Sea $A$ un anillo.
Las condiciones siguientes son equivalentes:
\begin{enumerate}[label=\emph{\alph*)}]
  \item $A$ es un anillo de Jacobson.
  \item Para todo ideal primo no maximal $\mathfrak{p}$ de $A$ y todo $f\neq 0$ en $B=A/\mathfrak{p}$, $B_f$ no es un cuerpo.
  \item[\rm b$'$)] Toda $A$-álgebra de tipo finito $K$ que sea un cuerpo es una $A$-álgebra finita.
\end{enumerate}
\end{proposition}

\begin{proof}
Sabemos que $a)$ implica $b')$ (Bourbaki, \emph{Alg.\ comm.}, cap.~V, \S\,3, n$^\circ$\,4, cor.~3 del th.~3).
Además, el núcleo del homomorfismo $A\to K$ es entonces un ideal \emph{maximal} $\mathfrak{m}$ de $A$, y $K$ es una extensión finita de $A/\mathfrak{m}$ (\emph{loc.\ cit.}).
Es trivial que $b')$ implica $b)$, pues $A/\mathfrak{p}=B$ no es un cuerpo, toda $B$-álgebra de tipo finito es una $A$-álgebra de tipo finito y $B_f$ es una $B$-álgebra de tipo finito.
Queda por mostrar que $b)$ implica $a)$.
Utilizaremos el lema siguiente:
\end{proof}

\begin{lemma}[10.4.5.1]
\label{IV.10.4.5.1}
Sea $X$ un espacio topológico que tenga la propiedad siguiente: para todo subconjunto localmente cerrado $Z\neq\varnothing$ de $X$, existe un subconjunto $Z'$ de $Z$, localmente cerrado en $X$ (o en $Z$, lo que equivale a lo mismo), y un punto $x\in Z'$, cerrado en $Z'$.
Entonces, para que $X$ sea un espacio de Jacobson, es necesario y suficiente que ningún punto no cerrado $x$ de $X$ sea aislado en $\overline{\{x\}}$.
\end{lemma}

\begin{proof}
Si $X$ es un espacio de Jacobson, un punto no cerrado $x$ de $X$ no puede ser aislado en $\overline{\{x\}}$, pues esto significaría que existe un abierto $U$ de $X$ tal que $U\cap\overline{\{x\}}=\{x\}$; pero $U\cap\overline{\{x\}}$ es localmente cerrado en $X$ y no contendría ningún punto cerrado en $X$, lo que contradice la hipótesis de que $X$ es un espacio de Jacobson.
Recíprocamente, supongamos satisfechas las condiciones del enunciado; entonces el conjunto $X_0$ de los puntos cerrados de $X$ coincide con el conjunto de los $x\in X$ que son aislados en $\overline{\{x\}}$; pero se deduce de \sref{IV.5.1.10.1} que
\oldpage[IV-3]{103}
este conjunto es \emph{muy denso} en $X$, luego $X$ es un espacio de Jacobson por definición.
\end{proof}

Recordemos \sref{IV.5.1.10} que la hipótesis formulada sobre $X$ en \sref{IV.10.4.5.1} siempre se satisface cuando el espacio subyacente a $X$ es un \emph{preesquema}.

Volviendo a la demostración de \sref{IV.10.4.5}, la condición $b)$ significa que, para todo punto no cerrado $x$ de $\operatorname{Spec}(A)$, el punto genérico $x$ de $\operatorname{Spec}(A/\mathfrak{j}_x)=V(\mathfrak{j}_x)=\overline{\{x\}}$ no es aislado en $\overline{\{x\}}$; por tanto, $b)$ implica $a)$ por el lema \sref{IV.10.4.5.1} y \sref{IV.10.4.4}.

\begin{corollary}[10.4.6]
\label{IV.10.4.6}
Toda álgebra de tipo finito $B$ sobre un anillo de Jacobson $A$ es un anillo de Jacobson, y la imagen inversa en $A$ de todo ideal maximal de $B$ es un ideal maximal de $A$.
En particular, toda álgebra de tipo finito sobre un cuerpo o sobre $\mathbf{Z}$ es un anillo de Jacobson.
\end{corollary}

\begin{proof}
Una $B$-álgebra de tipo finito $K$ que sea un cuerpo es también una $A$-álgebra de tipo finito, y por tanto un $A$-módulo de tipo finito y \emph{a fortiori} un $B$-módulo de tipo finito, de donde la primera afirmación; la segunda se probó en el curso de la demostración de \sref{IV.10.4.5}, aplicada a $K=B/\mathfrak{m}$.
\end{proof}

\begin{corollary}[10.4.7]
\label{IV.10.4.7}
Si $X$ es un preesquema de Jacobson y $f:Y\to X$ es un morfismo localmente de tipo finito, entonces $Y$ es un preesquema de Jacobson y la imagen por $f$ de todo punto cerrado de $Y$ es un punto cerrado de $X$.
\end{corollary}

\begin{proof}
Como la cuestión es local sobre $X$ y sobre $Y$, nos reducimos al caso en que $X$ e $Y$ son afines, y el corolario se deduce entonces de \sref{IV.10.4.6}.
\end{proof}

\begin{corollary}[10.4.8]
\label{IV.10.4.8}
Si $k$ es un cuerpo algebraicamente cerrado y $X$ un $k$-preesquema localmente de tipo finito, el conjunto de los puntos racionales de $X$ sobre $k$ es muy denso en $X$.
\end{corollary}

\begin{proof}
En efecto, $X$ es un preesquema de Jacobson \sref{IV.10.4.7} y los puntos cerrados son exactamente los puntos racionales sobre $k$ \sref[I]{I.6.4.2}.
\end{proof}

\begin{env}[10.4.9]
\label{IV.10.4.9}
El hecho de que los preesquemas localmente de tipo finito sobre un cuerpo o sobre $\mathbf{Z}$ sean preesquemas de Jacobson es particularmente importante, por la posibilidad de reducirse a este caso en muchas cuestiones de geometría algebraica \sref{IV.8.1.2},~$c)$.
Daremos dos ejemplos:
\end{env}

\begin{env}[10.4.10]
\label{IV.10.4.10}
\textbf{Aplicaciones (10.4.10) :~I~: Demostración de \sref{IV.6.15.9}.}
Sea $k$ un cuerpo separablemente cerrado y $X$ un $k$-preesquema localmente de tipo finito sobre $k$ y \emph{unirramificado}.
Sabemos que el cierre entero de una $k$-álgebra íntegra de tipo finito $A$ en una extensión finita de su cuerpo de fracciones es una $A$-álgebra finita (Bourbaki, \emph{Alg.\ comm.}, cap.~V, \S\,3, n$^\circ$\,2, th.~2), luego toda $k$-álgebra de tipo finito es un anillo universalmente japonés; concluimos que el conjunto de los puntos $x\in X$ en los que $X$ es geométricamente unirramificado es localmente construible \sref{IV.9.7.10}.
Pero la hipótesis y el lema \sref{IV.6.15.8} implican que este conjunto contiene todos los puntos \emph{cerrados} de $X$.
La conclusión se deduce por tanto de \sref{IV.10.4.6}, \sref{IV.10.3.1} y de la biyectividad de la aplicación canónica $\mathfrak{Lc}(X)\to\mathfrak{Lc}(X_0)$ (donde $X_0$ es el conjunto de los puntos cerrados de $X$) \sref{IV.10.2.7}.
\end{env}

\begin{env}[10.4.11]
\label{IV.10.4.11}
\oldpage[IV-3]{104}
\textbf{Aplicaciones :~II~: Proposición (10.4.11).}
Sea $S$ un preesquema y $X$ un $S$-preesquema de tipo finito.
Todo $S$-endomorfismo de $X$ que sea radicial es sobreyectivo (y por tanto biyectivo).
\end{env}

\begin{proof}
Sea $f:X\to S$ el morfismo estructural, $g:X\to X$ el $S$-endomorfismo considerado y, para todo $s\in S$, sea $g_s$ el morfismo deducido de $g$ por el cambio de base $\operatorname{Spec}(k(s))\to S$, que es un $k(s)$-endomorfismo de la fibra $f^{-1}(s)=X\times_S\operatorname{Spec}(k(s))$.

Para probar que $g$ es sobreyectivo, basta probar que $g_s$ es sobreyectivo para todo $s\in S$; podemos por tanto (en virtud de \sref[I]{I.3.5.7}) reducirnos al caso en que $S=\operatorname{Spec}(k)$ es el espectro de un cuerpo $k$, en cuyo caso $f$ es un morfismo de presentación finita, pues $S$ es noetheriano.
Aplicando \sref{IV.8.9.1} y \sref{IV.8.10.5},~(vii), nos reducimos al caso en que $S=\operatorname{Spec}(A)$, donde $A$ es una sub-$\mathbf{Z}$-álgebra de tipo finito de $k$.
Pero entonces $X$ es un preesquema de Jacobson \sref{IV.10.4.7}, y $g(X)$ es construible en $X$ \sref[I]{I.1.8.5}; por tanto, para probar que $g(X)=X$, basta mostrar que $g(X)$ contiene todos los puntos \emph{cerrados} de $X$ \sref{IV.10.3.1}.

\begin{lemma}[10.4.11.1]
\label{IV.10.4.11.1}
Sea $Y$ un $\mathbf{Z}$-preesquema de tipo finito.
\begin{enumerate}[label=(\roman*)]
  \item Para que un punto $y\in Y$ sea cerrado, es necesario y suficiente que $k(y)$ sea un cuerpo finito.
  \item Para todo número primo $p$ y todo entero $d\geq 1$, el conjunto de los puntos $y\in Y$ tales que $k(y)$ es una extensión de $\mathbf{F}_p$ de grado divisor de $d$ es finito.
\end{enumerate}
\end{lemma}

La afirmación (i) se deduce del hecho de que la imagen de un punto cerrado $y\in Y$ en $\operatorname{Spec}(\mathbf{Z})$ es un punto cerrado \sref{IV.10.4.7}, es decir, un número primo, y de \sref[I]{I.6.4.2}.
Por otra parte, como $Y$ es una unión finita de abiertos afines de tipo finito sobre $\mathbf{Z}$, podemos limitarnos a probar (ii) en el caso en que $Y=\operatorname{Spec}(C)$, donde $C$ es una $\mathbf{Z}$-álgebra de tipo finito.
Ahora bien, los puntos $y\in Y$ tales que el grado de $k(y)$ sobre $\mathbf{F}_p$ divide a $d$ son exactamente los núcleos de los homomorfismos $C\to\mathbf{F}_{p^d}$; pero, si $(t_i)$ es un sistema de generadores de la $\mathbf{Z}$-álgebra $C$, todo homomorfismo de $C$ queda determinado por sus valores sobre los elementos $t_i$, y por tanto solo hay un número finito de homomorfismos de $C$ en un cuerpo finito.

Siendo esto así, como $X$ es un $\mathbf{Z}$-preesquema de tipo finito, sea $T_{p,d}$ el conjunto de los puntos cerrados $z\in X$ tales que $k(z)$ es una extensión de $\mathbf{F}_p$ de grado divisor de $d$; se deduce de \sref{IV.10.4.11.1} que el conjunto $T_{p,d}$ es finito y que el conjunto de los puntos cerrados de $X$ es la unión de los $T_{p,d}$.
Además, si $z\in T_{p,d}$ y $h$ es un endomorfismo cualquiera de $X$, $k(h(z))$ es isomorfo a un subcuerpo de $k(z)$, luego $h(z)\in T_{p,d}$; dicho de otro modo, $T_{p,d}$ es \emph{estable} bajo todo endomorfismo de $X$.
Como $g$ es inyectivo por hipótesis, su restricción a $T_{p,d}$ es \emph{biyectiva}, pues $T_{p,d}$ es finito, lo que termina de probar la proposición.

Veremos más adelante \sref{IV.17.9.6} que, si además se supone que $X$ es un $S$-preesquema de presentación finita y, por otra parte, que $g$ es un monomorfismo, puede afirmarse que $g$ es un \emph{automorfismo} de $X$.
\end{proof}

\subsection{Preesquemas de Jacobson noetherianos}
\label{subsection:IV.10.5}

\begin{proposition}[10.5.1]
\label{IV.10.5.1}
Sea $B$ un anillo noetheriano íntegro.
Las condiciones equivalentes $a)$ a $d)$ de \sref{IV.10.4.4} son entonces también equivalentes a las siguientes:
\begin{enumerate}[label=\emph{\alph*)}]
  \setcounter{enumi}{4}
  \item $\operatorname{Spec}(B)$ es finito.
  \item $B$ es un anillo semilocal de dimensión $\leq 1$.
\end{enumerate}
\end{proposition}

\begin{proof}
\oldpage[IV-3]{105}
En efecto, se deduce del teorema de Artin--Tate \sref[0]{0.16.3.3} que las condiciones $a)$ y $f)$ son equivalentes.
La condición $f)$ implica que $\operatorname{Spec}(B)$ es la unión del conjunto de sus puntos cerrados y de su punto genérico, luego $f)$ implica $e)$ sin suponer que $B$ sea noetheriano.
Finalmente, $e)$ implica $d)$ sin suponer que $B$ sea noetheriano, pues el punto genérico $x$ de $X$ es el complementario de la unión de las adherencias $\overline{\{y\}}$, donde $y$ recorre el conjunto de los puntos $y\neq x$, y como estos puntos son finitos en número, $\{x\}$ es el complementario de un cerrado de $X$.
\end{proof}

\begin{corollary}[10.5.2]
\label{IV.10.5.2}
Para que un anillo noetheriano $A$ sea un anillo de Jacobson, es necesario y suficiente que no exista ningún ideal primo $\mathfrak{p}$ de $A$ tal que $A/\mathfrak{p}$ sea un anillo semilocal de dimensión $1$.
\end{corollary}

\begin{proof}
Esto se deduce inmediatamente de \sref{IV.10.5.1} y de la condición $b)$ de \sref{IV.10.4.5}, pues los ideales primos $\mathfrak{p}$ de $A$ tales que $A/\mathfrak{p}$ es semilocal de dimensión~$0$ son los ideales maximales de $A$.
\end{proof}

\begin{corollary}[10.5.3]
\label{IV.10.5.3}
Sea $X$ un preesquema noetheriano irreducible.
Las condiciones siguientes son equivalentes:
\begin{enumerate}[label=\emph{\alph*)}]
  \item El punto genérico de $X$ es aislado.
  \item $X$ es finito.
  \item $X$ es de dimensión $\leq 1$ y el conjunto de sus puntos cerrados es finito.
\end{enumerate}
\end{corollary}

\begin{proof}
Existe por hipótesis un número finito de abiertos afines irreducibles $U_i$ ($1\leq i\leq n$) que recubren $X$ y cada uno de los cuales contiene el punto genérico de $X$; basta probar la equivalencia de $a)$, $b)$ y $c)$ para cada uno de los $(U_i)_{\mathrm{red}}$ (teniendo en cuenta \sref[0]{0.14.1.7}).
Pero esta equivalencia se deduce de \sref{IV.10.5.1}.
\end{proof}

\begin{remark}[10.5.4]
\label{IV.10.5.4}
Un preesquema noetheriano $X$ que satisfaga las condiciones equivalentes de \sref{IV.10.5.3} no es necesariamente un esquema afín; de hecho, puede incluso ser \emph{no separado}, como se ve sustituyendo, en el ejemplo \sref[I]{I.5.5.11} de la ``recta afín con punto duplicado'', $X_1$ y $X_2$ por el espectro del anillo de valoración discreta $(K[s])_{(s)}$, y $U_{12}$ y $U_{21}$ por el abierto reducido al punto genérico en $X_1$ y $X_2$ respectivamente; el preesquema no separado $X$ así obtenido tiene exactamente $3$ puntos.
\end{remark}

\begin{proposition}[10.5.5]
\label{IV.10.5.5}
Sea $X$ un preesquema noetheriano.
\begin{enumerate}[label=(\roman*)]
  \item El conjunto $X_0$ de los puntos $x\in X$ tales que $\overline{\{x\}}$ es finito es muy denso en $X$.
  \item Para que $X$ sea un preesquema de Jacobson, es necesario y suficiente que no exista ningún subpreesquema de $X$ isomorfo al espectro de un anillo semilocal íntegro de dimensión $1$.
\end{enumerate}
\end{proposition}

\begin{proof}
La condición de que $\overline{\{x\}}$ sea finito equivale aquí, por \sref{IV.10.5.3}, a que $x$ sea aislado en $\overline{\{x\}}$, pues este es el espacio subyacente a un subpreesquema (noetheriano) de $X$; la afirmación (i) se deduce por tanto de \sref{IV.5.1.10.1}.
Análogamente, teniendo en cuenta \sref{IV.10.4.5.1}, para mostrar la afirmación (ii), observemos que, para que un punto no cerrado $x$ de $X$ pertenezca a $X_0$, es necesario y suficiente que el subpreesquema cerrado íntegro $Y$ de $X$ cuyo espacio subyacente es $\overline{\{x\}}$ sea de dimensión $1$ y finito, y por tanto una unión finita de subpreesquemas $U_i$ abiertos (en $Y$) y afines que sean espectros de anillos semilocales íntegros de dimensión $1$.
Recíprocamente, si existe un subpreesquema $Z$ de $X$ que sea el espectro de un anillo semilocal íntegro de dimensión $1$, $Z$ no es un preesquema de Jacobson \sref{IV.10.5.2}, y por tanto tampoco lo es $X$ \sref{IV.10.3.2}.
\end{proof}

\oldpage[IV-3]{106}

\begin{remark}[10.5.6]
\label{IV.10.5.6}
La afirmación (ii) de \sref{IV.10.5.5} sigue siendo válida cuando $X$ es \emph{localmente noetheriano}: en efecto, si $(V_\alpha)$ es un recubrimiento de $X$ por abiertos afines (noetherianos), todo subpreesquema de uno de los $V_\alpha$ es un subpreesquema de $X$; recíprocamente, si un subpreesquema $Z$ de $X$ es isomorfo al espectro de un anillo semilocal íntegro de dimensión $1$, existe un $\alpha$ tal que $V_\alpha\cap Z$ contiene un abierto afín $U$ de $Z$ que no se reduce al punto genérico de $Z$, y que por tanto también es el espectro de un anillo semilocal íntegro de dimensión $1$.
Concluimos por \sref{IV.10.3.3}.
\end{remark}

\begin{proposition}[10.5.7]
\label{IV.10.5.7}
Sea $X$ un preesquema localmente noetheriano y $Y$ un cerrado de $X$ tal que todo cerrado no vacío de $X$ corta a $Y$.
Entonces el subpreesquema inducido sobre el abierto $X-Y$ es un preesquema de Jacobson.
\end{proposition}

\begin{proof}
Apliquemos el criterio \sref{IV.10.5.5},~(ii), y supongamos que existe en $X-Y$ un subpreesquema $Z$ que es el espectro de un anillo semilocal íntegro de dimensión $1$. El punto genérico $z$ de $Z$ es entonces aislado en $Z$ (o en la adherencia $\overline{Z}$ de $Z$ en $X$), y $Z$ es distinto de $\{z\}$.
Sea $y\neq z$ un punto de $Z$; como no pertenece a $Y$, no es un punto cerrado de $X$, y su adherencia $\overline{\{y\}}$ en $X$ corta a $Y$ en un punto $x\neq y$, que es por tanto una especialización de $y$.
La existencia de la cadena $\{x\}\subset\overline{\{y\}}\subset\overline{\{z\}}$ muestra entonces que la dimensión de $\overline{\{z\}}$ sería $\geq 2$, y lo mismo ocurriría con la dimensión de $U\cap\overline{\{z\}}$, donde $U$ es un entorno abierto afín (y por tanto noetheriano) de $x$ en $X$.
Pero esto contradice que el punto genérico $z$ de $U\cap\overline{\{z\}}$ sea aislado en $U\cap\overline{\{z\}}$ \sref{IV.10.5.3}.
\end{proof}

\begin{corollary}[10.5.8]
\label{IV.10.5.8}
Sea $A$ un anillo noetheriano; para todo elemento $f$ del radical $\mathfrak{R}$ de $A$, el anillo $A_f$ es un anillo de Jacobson, y el abierto $\operatorname{Spec}(A)-V(\mathfrak{R})$ es un esquema de Jacobson.
\end{corollary}

\begin{proof}
Si $X=\operatorname{Spec}(A)$ y $Y=V(\mathfrak{J})$, donde $\mathfrak{J}$ es un ideal de $A$, decir que todo cerrado no vacío de $X$ corta a $Y$ equivale \sref[0]{0.2.1.3} a decir que $Y$ contiene todos los \emph{puntos cerrados} de $X$, o equivalentemente que $\mathfrak{J}$ está \emph{contenido en el radical} $\mathfrak{R}$ de $A$.
Si $f\in\mathfrak{R}$, el abierto $D(f)$ no corta a $V(\mathfrak{R})$, y es un espacio de Jacobson en virtud de \sref{IV.10.5.7}.
\end{proof}

\begin{corollary}[10.5.9]
\label{IV.10.5.9}
Sea $A$ un anillo local noetheriano, $\mathfrak{m}$ su ideal maximal, $X=\operatorname{Spec}(A)$ y $Y=X-\{\mathfrak{m}\}$; entonces $Y$ es un esquema de Jacobson, cuyos puntos cerrados son los ideales primos $\mathfrak{p}\in\operatorname{Spec}(A)$ tales que $\dim(A/\mathfrak{p})=1$.
\end{corollary}

\begin{proof}
La primera afirmación es un caso particular de \sref{IV.10.5.8}; por otra parte, los puntos cerrados de $Y$ son los ideales primos $\mathfrak{p}$ de $A$ que son elementos maximales del conjunto de los ideales primos $\neq\mathfrak{m}$, lo que, por definición de la dimensión, significa que $\dim(A/\mathfrak{p})=1$.
\end{proof}

\begin{proposition}[10.5.10]
\label{IV.10.5.10}
Sea $A$ un anillo local noetheriano, reducido, completo y que no sea un cuerpo.
Entonces las intersecciones finitas de los núcleos de los homomorfismos locales de $A$ en anillos de valoración discreta $V$, que hacen de $V$ una $A$-álgebra finita, forman una base de filtro convergente a $0$ para la topología ádica de $A$.
\end{proposition}

\begin{proof}
Basta (Bourbaki, \emph{Alg.\ comm.}, cap.~III, \S\,2, n$^\circ$\,7, prop.~8) probar que la intersección de los núcleos considerados en el enunciado se reduce a $0$.
Supongamos primero que $A$ es íntegro y de dimensión $1$; en virtud del teorema de Nagata \sref[0]{0.23.1.5} y \sref[0]{0.23.1.6}, el cierre entero $A'$ de $A$ es un anillo local íntegro completo,
\oldpage[IV-3]{107}
íntegramente cerrado y de dimensión $1$, y una $A$-álgebra finita; es por tanto un anillo de valoración discreta, y la proposición se deduce inmediatamente en este caso.

Pasemos al caso general; sea $x$ el punto cerrado de $X=\operatorname{Spec}(A)$ y pongamos $Y=X-\{x\}$; sabemos \sref{IV.10.5.9} que $Y$ es un preesquema de Jacobson.
Mostremos que la intersección de los ideales primos $\mathfrak{p}$ de $A$ tales que $\dim(A/\mathfrak{p})=1$ se reduce a $0$; la proposición se seguirá entonces, pues, para cada uno de estos ideales $\mathfrak{p}$, la intersección de los núcleos de los homomorfismos locales de $A/\mathfrak{p}$ en anillos de valoración discreta $V$ que hacen de $V$ una $(A/\mathfrak{p})$-álgebra finita se reduce a $0$.
Pero decir que la intersección de estos ideales primos se reduce a $0$ significa que el conjunto de dichos ideales es denso en $X$, o equivalentemente en $Y$ (puesto que $Y$ es denso en $X$), y esto se deduce inmediatamente de \sref{IV.10.5.9}.
\end{proof}

\subsection{Dimensión en los preesquemas de Jacobson}
\label{subsection:IV.10.6}

Los resultados de esta sección precisan, en ciertos casos, y generalizan resultados del \S\,5.

\begin{proposition}[10.6.1]
\label{IV.10.6.1}
Sea $S$ un preesquema localmente noetheriano que satisfaga además las condiciones: $1^\circ$~$S$ es un preesquema de Jacobson; $2^\circ$~para todo $s\in S$, $\sh{O}_{S,s}$ es universalmente catenario \sref{IV.5.6.2}; $3^\circ$~toda componente irreducible $S'$ de $S$ es equicodimensional (dicho de otro modo, para todo punto cerrado $s$ de $S'$ y todo subpreesquema de $S$ cuyo espacio subyacente sea $S'$, se tiene $\dim(S')=\dim(\sh{O}_{S',s})$).
Se tienen entonces las propiedades siguientes:
\begin{enumerate}[label=(\roman*)]
  \item Para todo morfismo localmente de tipo finito $g:X\to S$, $X$ satisface las condiciones $1^\circ$, $2^\circ$ y $3^\circ$ anteriores.
En particular, si $X$ es equidimensional (por ejemplo, si $X$ es irreducible), $X$ es biequidimensional (dicho de otro modo, $X$ es catenario y para todo punto cerrado $x$ de $X$ se tiene $\dim(\sh{O}_{X,x})=\dim(X)$ \sref[0]{0.14.3.3}).
  \item Sean $X$, $Y$ dos $S$-preesquemas localmente de tipo finito sobre $S$ y $f:X\to Y$ un $S$-morfismo; supongamos que $X$ es irreducible y que $f$ es dominante.
Si $\xi$ (resp.\ $\eta$) es el punto genérico de $X$ (resp.\ $Y$) y $e=\dim(f^{-1}(\eta))=\deg.\mathrm{tr}_{k(\eta)}k(\xi)$, se tiene
\begin{equation}
\label{IV.10.6.1.1}
\tag{10.6.1.1}
\dim(X)=\dim(Y)+e.
\end{equation}
  \item Sean $X$, $Y$ dos $S$-preesquemas localmente de tipo finito sobre $S$, $f:X\to Y$ un $S$-morfismo y $n$ un entero no negativo tal que $\dim(f^{-1}(y))\geq n$ (resp.\ $\dim(f^{-1}(y))\leq n$) para todo $y\in Y$.
Entonces se tiene
\begin{equation}
\label{IV.10.6.1.2}
\tag{10.6.1.2}
\dim(X)\geq\dim(Y)+n
\end{equation}
(resp.
\begin{equation}
\label{IV.10.6.1.3}
\tag{10.6.1.3}
\dim(X)\leq\dim(Y)+n).
\end{equation}
\end{enumerate}
\end{proposition}

\begin{proof}
(i) La propiedad $1^\circ$ para $X$ se deduce de \sref{IV.10.4.7}.
Para todo $x\in X$, $\sh{O}_{X,x}$ es el anillo local, en un ideal primo, de una $\sh{O}_{S,g(x)}$-álgebra de tipo finito, y el homomorfismo $\sh{O}_{S,g(x)}\to\sh{O}_{X,x}$
\oldpage[IV-3]{108}
es local; por tanto \sref{IV.5.6.3},~(iv), $\sh{O}_{X,x}$ es universalmente catenario.
Para mostrar que $X$ satisface la condición $3^\circ$, consideremos varios casos:

\emph{a)} $X$ es un subpreesquema cerrado irreducible de $S$; sea $S'$ una componente irreducible de $S$ que contiene a $X$, $\xi$ el punto genérico de $X$ y $x$ un punto cerrado de $X$; para todo $s\in S'$, $\sh{O}_{S',s}$, cociente de $\sh{O}_{S,s}$, es catenario \sref{IV.5.6.1}, luego las condiciones $2^\circ$ y $3^\circ$ implican que $S'$ es biequidimensional \sref[0]{0.14.3.3}.
En virtud de \sref{IV.5.1.2} y de \sref[0]{0.14.3.3.2}, se tiene por tanto
\[
  \dim(\sh{O}_{X,x})=\dim(\sh{O}_{S',x})-\dim(\sh{O}_{S',\xi})=\dim(S')-\dim(\sh{O}_{S',\xi})
\]
lo que prueba que $\dim(\sh{O}_{X,x})$ no depende del punto cerrado $x$ considerado, de donde la afirmación en este caso \sref{IV.5.1.4}.

\emph{b)} $X$ es irreducible y $g$ dominante.
Entonces, para todo punto cerrado $x$ de $X$, $g(x)$ es cerrado en $S$ por \sref{IV.10.4.7}; como $\sh{O}_{S,g(x)}$ es universalmente catenario, se deduce de \sref{IV.5.6.5.3} que se tiene
\[
  \dim(\sh{O}_{X,x})=\dim(\sh{O}_{S,g(x)})+e
\]
donde $e=\dim(g^{-1}(\zeta))$, siendo $\zeta$ el punto genérico de $S$.
Como $\dim(S)=\dim(\sh{O}_{S,g(x)})$ en virtud de la condición $3^\circ$ para $S$, se tiene $\dim(\sh{O}_{X,x})=\dim(S)+e$ para todo punto cerrado $x\in X$; esto prueba la condición $3^\circ$ para $X$ \sref{IV.5.1.4}, y al mismo tiempo la fórmula
\begin{equation}
\label{IV.10.6.1.4}
\tag{10.6.1.4}
\dim(X)=\dim(S)+e.
\end{equation}

\emph{c) Caso general.}
Considerando un subpreesquema reducido $X'$ de $X$ cuyo espacio subyacente sea una componente irreducible de $X$, nos reducimos al caso en que $X$ es íntegro; utilizando \sref[I]{I.5.2.2} y el caso $a)$ probado anteriormente, podemos sustituir entonces $S$ por el subpreesquema reducido cuyo espacio subyacente es $\overline{g(X)}$; quedamos así reducidos al caso $b)$, lo que termina de probar (i).

(ii) Como el morfismo $f$ es localmente de tipo finito \sref[I]{I.1.3.4}, podemos aplicar los resultados de (i) sustituyendo $S$ por $Y$; además, como $X$ es irreducible y $f$ dominante, también podemos sustituir $S$ por $Y$ en \sref{IV.10.6.1.4}, lo que da \sref{IV.10.6.1.1}.

(iii) La afirmación relativa al caso en que $\dim(f^{-1}(y))\leq n$ para todo $y\in Y$ ya se ha probado bajo hipótesis más generales en \sref{IV.5.6.7}.
Supongamos que $\dim(f^{-1}(y))\geq n$ para todo $y\in Y$, y consideremos un punto genérico $\eta$ de una componente irreducible $Y'$ de $Y$; existe al menos una componente irreducible $Z$ de $f^{-1}(\eta)$ de dimensión $\geq n$; si $\xi$ es el punto genérico de $Z$, $\xi$ también es un punto genérico de una componente irreducible $X'$ de $X$ tal que $f(X')$ es denso en $Y'$ \sref[0]{0.2.1.8}.
Consideremos los subpreesquemas reducidos de $X$, $Y$ que tienen respectivamente por espacios subyacentes $X'$, $Y'$, y la restricción $X'\to Y'$ de $f$ \sref[I]{I.5.2.2}; se deduce entonces de (ii) que $\dim(X')\geq\dim(Y')+n$, y \emph{a fortiori} $\dim(X)\geq\dim(Y')+n$; como esto es cierto para toda componente irreducible $Y'$ de $Y$, concluimos la desigualdad \sref{IV.10.6.1.2}.
\end{proof}

\begin{corollary}[10.6.2]
\label{IV.10.6.2}
Supongamos que $X$ satisface las condiciones $1^\circ$, $2^\circ$ y $3^\circ$ de \sref{IV.10.6.1}.
Entonces, para todo abierto denso $U$ de $X$, se tiene $\dim(U)=\dim(X)$.
\end{corollary}

\begin{proof}
\oldpage[IV-3]{109}
Sabemos en efecto que $\dim(U)\leq\dim(X)$ \sref[0]{0.14.1.4}; además, como $X$ es un preesquema de Jacobson, $U$ contiene un punto cerrado de cada componente irreducible de $X$, luego $\dim(U)=\dim(X)$ en virtud de \sref{IV.10.6.1} y \sref[0]{0.14.1.2.1}.
\end{proof}

\begin{proposition}[10.6.3]
\label{IV.10.6.3}
Supongamos que $X$ satisface las condiciones $1^\circ$, $2^\circ$ y $3^\circ$ de \sref{IV.10.6.1}, y sea $Y$ un cerrado de $X$.
Para todo $x\in Y$ y todo entorno abierto $U$ de $x$ en $X$ que no corte las componentes irreducibles de $Y$ que no contienen a $x$, se tiene
\begin{equation}
\label{IV.10.6.3.1}
\tag{10.6.3.1}
\dim_x(Y)=\dim(U\cap Y)=\dim\overline{\{x\}}+\operatorname{codim}(\overline{\{x\}},Y)=\sup_i\dim(Y_i)
\end{equation}
donde los $Y_i$ ($1\leq i\leq m$) son las componentes irreducibles de $Y$ que contienen a $x$.
\end{proposition}

\begin{proof}
Considerando un subpreesquema cerrado de $X$ cuyo espacio subyacente sea $Y$, podemos limitarnos, en virtud de \sref{IV.10.6.1},~(i), al caso en que $Y=X$.
Por la elección de $U$, $U$ es la unión de los $U\cap X_i$, luego $\dim(U)=\sup_i\dim(U\cap X_i)$; pero, por \sref{IV.10.6.2} aplicado a un subpreesquema cerrado de $X$ cuyo espacio subyacente sea $X_i$, se tiene (teniendo en cuenta \sref{IV.10.6.1},~(i)) $\dim(U\cap X_i)=\dim(X_i)$; esto prueba que el segundo y el cuarto miembros de \sref{IV.10.6.3.1} son iguales.
Como $U\cap X_i$ es biequidimensional \sref{IV.10.6.1},~(i) (pues la inmersión $U\cap X_i\to X_i$ es de tipo finito \sref[I]{I.6.3.5}), se tiene

\begin{equation}
\label{IV.10.6.3.2}
\tag{10.6.3.2}
\dim(X_i)=\dim(U\cap X_i)=\dim(U\cap\overline{\{x\}})+\operatorname{codim}(U\cap\overline{\{x\}},U\cap X_i)
\end{equation}
\sref[0]{0.14.3.5}.
Por \sref{IV.10.6.2}, aplicado a un subpreesquema cerrado de $X$ cuyo espacio subyacente sea $\overline{\{x\}}$, $\dim(U\cap\overline{\{x\}})=\dim(\overline{\{x\}})$; como, por las mismas razones, también se tiene

\begin{equation}
\label{IV.10.6.3.3}
\tag{10.6.3.3}
\dim(X_i)=\dim(\overline{\{x\}})+\operatorname{codim}(\overline{\{x\}},X_i)
\end{equation}
se obtiene
\[
  \dim(U)=\dim(\overline{\{x\}})+\sup_i\operatorname{codim}(\overline{\{x\}},X_i)=\dim(\overline{\{x\}})+\operatorname{codim}(\overline{\{x\}},X)
\]
por definición de la codimensión \sref[0]{0.14.2.1}.
Esto muestra que $\dim(U)$ es \emph{independiente} del entorno abierto $U$ de $x$ que satisface las condiciones del enunciado, y por tanto es igual a $\dim_x(X)$, por \sref[0]{0.14.1.4.1}.
\end{proof}

\begin{corollary}[10.6.4]
\label{IV.10.6.4}
Bajo las hipótesis de \sref{IV.10.6.3}, sea $\sh{F}$ un $\sh{O}_X$-módulo coherente y $Y$ su soporte.
Para todo $x\in Y$, se tiene
\begin{equation}
\label{IV.10.6.4.1}
\tag{10.6.4.1}
\dim(\overline{\{x\}})+\dim(\sh{F}_x)=\dim_x Y.
\end{equation}
\end{corollary}

\begin{proof}
En efecto, esto se deduce de \sref{IV.10.6.3.1} y de la fórmula $\dim(\sh{F}_x)=\operatorname{codim}(\overline{\{x\}},Y)$ \sref{IV.5.1.12.2}.
\end{proof}

\subsection{Ejemplos y contraejemplos}
\label{subsection:IV.10.7}

\begin{env}[10.7.1]
\label{IV.10.7.1}
Sea $S$ un preesquema localmente noetheriano de dimensión $\leq 1$ y supongamos que $S$ es un \emph{preesquema de Jacobson}; cuando $S$ es noetheriano, esto equivale a decir que las componentes irreducibles de $S$ de dimensión $1$ son infinitas, pues todo $x\in S$ que no sea un punto cerrado es el punto genérico de una componente tal \sref{IV.10.4.5} y \sref{IV.10.5.4}.
Entonces $S$ satisface también las condiciones $2^\circ$ y $3^\circ$ de \sref{IV.10.6.1}: en efecto, todo anillo local $\sh{O}_{S,s}$ es de dimensión $0$ o $1$, y por tanto universalmente catenario \sref{IV.7.2.9}; por otra parte, una componente irreducible $S'$ de $S$ o bien se reduce
\oldpage[IV-3]{110}
a un punto o bien es de dimensión $1$, y para todo punto cerrado $s\in S'$, $\sh{O}_{S',s}$ es necesariamente de dimensión $1$.

De estas observaciones y de \sref{IV.10.6.1} se deduce que todo preesquema localmente de tipo finito sobre $S$ satisface también las propiedades $1^\circ$, $2^\circ$ y $3^\circ$; este es en particular el caso de los preesquemas \emph{localmente de tipo finito sobre un cuerpo o sobre $\mathbf{Z}$}.
\end{env}

\begin{env}[10.7.2]
\label{IV.10.7.2}
Sea $A$ un anillo local noetheriano \emph{universalmente catenario} y sea $S$ el complementario del \emph{punto cerrado} $a$ en $X=\operatorname{Spec}(A)$.
Entonces $S$ satisface las condiciones $1^\circ$, $2^\circ$ y $3^\circ$ de \sref{IV.10.6.1}: en efecto, ya se ha visto que $S$ es un preesquema de Jacobson \sref{IV.10.5.9}; como $A$ es universalmente catenario, también lo son todos los anillos locales $A_\mathfrak{p}$ en los ideales primos de $A$ \sref{IV.5.6.3}.
Por otra parte, una componente irreducible $S'$ de $S$ es el complementario de $a$ en una componente irreducible $X'$ de $X$. Para todo punto cerrado $x$ de $S$, la adherencia de $x$ en $X$ es $\{x,a\}$, luego $\dim(\overline{\{x\}})=1$. Como $X'$ es biequidimensional por hipótesis \sref[0]{0.14.3.3}, se deduce que $\dim(\sh{O}_{X',x})=\dim(X')-1$ \sref[0]{0.14.3.2}, lo que prueba que $S'$ es equicodimensional \sref{IV.5.1.4}.
\end{env}

\begin{env}[10.7.3]
\label{IV.10.7.3}
Sea $A$ un anillo de valoración discreta; mostremos que en $B=A[T_1,\ldots,T_n]$, con $n$ indeterminadas, existen dos ideales maximales $\mathfrak{m}$, $\mathfrak{n}$ de \emph{alturas respectivas} $n$ y $n+1$.
Lo hemos visto en \sref{IV.5.2.5},~(i) para $n=1$; probémoslo por inducción sobre $n$.
Como $A[T_1,\ldots,T_{n+1}]$ es un $A[T_1,\ldots,T_n]$-módulo libre, y por tanto fielmente plano, existen ideales maximales $\mathfrak{m}'$ y $\mathfrak{n}'$ de $A[T_1,\ldots,T_{n+1}]$ situados respectivamente sobre $\mathfrak{m}$ y $\mathfrak{n}$ \sref[0]{0.6.5.1}; por \sref{IV.5.5.3}, estos ideales tienen alturas respectivas $n+1$ y $n+2$, lo que prueba la afirmación.
Supongamos $n\geq 2$ en lo que sigue.
Sea $\mathfrak{J}=\mathfrak{m}\cap\mathfrak{n}=\mathfrak{m}\mathfrak{n}$ y sea $R=1+\mathfrak{J}$, que es un subconjunto multiplicativo de $B$.
Si $B'=R^{-1}B$, el ideal $\mathfrak{J}'=R^{-1}\mathfrak{J}$ está contenido en el radical de $B'$ (Bourbaki, \emph{Alg\`ebre commutative}, cap.~III, \S3, n.~5, prop.~12).
Como espacio topológico, $X'=\operatorname{Spec}(B')$ se identifica con un subespacio de $X=\operatorname{Spec}(B)$, y en los puntos $x$ de $X'$ los anillos locales $\sh{O}_{X,x}$ y $\sh{O}_{X',x}$ son los mismos \sref[I]{I.1.6.2}.
Consideremos en $X'$ el cerrado $Y'=V(\mathfrak{J}')$, y pongamos $S=X'\setminus Y'$.
Entonces $S$ es un preesquema de Jacobson \sref{IV.10.5.7}, y es evidentemente irreducible y noetheriano; además, para todo $s\in S$, el anillo local $\sh{O}_{S,s}=\sh{O}_{X,s}$ es universalmente catenario por \sref{IV.5.6.3}, pues $A$ es universalmente catenario \sref{IV.5.6.4}.
Sin embargo, existen dos puntos cerrados $a,b$ de $S$ tales que $\sh{O}_{S,a}$ y $\sh{O}_{S,b}$ \emph{no tienen la misma dimensión}; dicho de otro modo, $S$ no satisface la condición $3^\circ$ de \sref{IV.10.6.1}.
Para verlo, consideremos los dos ideales maximales $\mathfrak{m}'=R^{-1}\mathfrak{m}$ y $\mathfrak{n}'=R^{-1}\mathfrak{n}$ de $B'$, que tienen alturas $n$ y $n+1$, respectivamente.
Se tiene $\mathfrak{J}'=\mathfrak{m}'\cap\mathfrak{n}'$, y por tanto $\mathfrak{J}'$ no está contenido en ningún ideal primo de $B'$ distinto de $\mathfrak{m}'$ y $\mathfrak{n}'$; estos son, en consecuencia, los únicos ideales maximales de $B'$.
Sean $a',b'$ los únicos puntos cerrados de $X'$, correspondientes a $\mathfrak{m}'$ y $\mathfrak{n}'$.
Existe en $B'$ un ideal primo no maximal $\mathfrak{p}'\subset\mathfrak{m}'$ que no está contenido en $\mathfrak{n}'$: basta mostrar que en $B$ existe un ideal primo no maximal contenido en $\mathfrak{m}$ y no en $\mathfrak{n}$; para ello puede, por ejemplo, considerarse la fibra de $\mathfrak{m}$ para el morfismo correspondiente a la inyección $A[T_1]\longrightarrow B=A[T_1,\ldots,T_n]$ y aplicar \sref{IV.6.1.2}.
Considerando una cadena maximal de ideales primos entre $\mathfrak{p}'$ y $\mathfrak{m}'$, y sustituyendo $\mathfrak{p}'$ por el penúltimo ideal de esta cadena, podemos por tanto suponer que el punto $a$ de $X'$ correspondiente a $\mathfrak{p}'$ tiene por adherencia $\{a,a'\}$ en $X'$.
Como $B'_{\mathfrak{m}'}$ es biequidimensional, $\mathfrak{p}'$ tiene entonces altura $n-1$.
Del mismo modo se construye un ideal primo no maximal $\mathfrak{q}'$ de $B'$ de altura $n$ tal que, si $b$ es el punto correspondiente de $X'$, la adherencia de $b$ en $X'$ es $\{b,b'\}$.
Así, $a$ y $b$ pertenecen a $S$ y son cerrados en $S$, y proporcionan por tanto el ejemplo buscado.
\end{env}

\subsection{Profundidad rectificada}
\label{subsection:IV.10.8}

\begin{definition}[10.8.1]
\label{IV.10.8.1}
Sea $X$ un preesquema localmente noetheriano y $\sh{F}$ un $\sh{O}_X$-Módulo coherente.
Para todo $x\in X$, llamamos \emph{profundidad rectificada} de $\sh{F}$ en el punto $x$ y denotamos $\operatorname{prof}^*_x(\sh{F})$ el número (entero $\geq 0$ o $+\infty$) igual a
\begin{equation}
\label{IV.10.8.1.1}
\tag{10.8.1.1}
\operatorname{prof}^*_x(\sh{F})=\operatorname{prof}(\sh{F}_x)+\dim(\overline{\{x\}})
\end{equation}
donde $\overline{\{x\}}$ es la adherencia del punto $x$ en $X$.
Para todo subconjunto $Z$ de $X$, llamamos \emph{profundidad rectificada de $\sh{F}$ a lo largo de $Z$} y denotamos $\operatorname{prof}^*_Z(\sh{F})$ el número
\begin{equation}
\label{IV.10.8.1.2}
\tag{10.8.1.2}
\operatorname{prof}^*_Z(\sh{F})=\inf_{x\in Z}\operatorname{prof}^*_x(\sh{F}).
\end{equation}
\end{definition}

Dicho de otro modo, para todo entero $n$, la relación $\operatorname{prof}^*_Z(\sh{F})\geq n$ equivale a $\operatorname{prof}^*_x(\sh{F})\geq n$ para todo $x\in Z$.
Si $Z=X$, se escribe $\operatorname{prof}^*(\sh{F})$ en lugar de $\operatorname{prof}^*_X(\sh{F})$.

\oldpage[IV-3]{111}

\begin{remarks}[10.8.2]
\label{IV.10.8.2}
(i) En todo punto cerrado $x\in X$, la profundidad rectificada es igual a la profundidad.

(ii) Sea $Y$ un subpreesquema cerrado de $X$, $j:Y\to X$ la inyección canónica y $\sh{G}$ un $\sh{O}_Y$-módulo coherente.
Sabemos \sref{IV.5.7.3},~(vi) que $\operatorname{prof}(\sh{G}_x)=\operatorname{prof}(j_*(\sh{G})_x)$ para todo $x\in Y$; deducimos por tanto que $\operatorname{prof}^*_x(\sh{G})=\operatorname{prof}^*_x(j_*(\sh{G}))$ para todo $x\in Y$.

(iii) La noción de profundidad rectificada solo presenta interés cuando es de carácter \emph{local}, es decir, cuando no cambia al sustituir $X$ por un entorno abierto arbitrario $U$ de $x$.
Esto exige evidentemente que $x$ no sea aislado en $\overline{\{x\}}$ cuando $x$ no es un punto cerrado de $X$, y por consiguiente que $X$ sea un \emph{preesquema de Jacobson} \sref{IV.10.4.5.1}; la mayoría de las veces, también será necesario saber que $\dim(U)=\dim(X)$ para todo abierto denso $U$ de $X$, y habrá que suponer por tanto que $X$ satisface además las condiciones $2^\circ$ y $3^\circ$ de \sref{IV.10.6.1}.
\end{remarks}

\begin{lemma}[10.8.3]
\label{IV.10.8.3}
Sea $X$ un preesquema regular y biequidimensional y $\sh{F}$ un $\sh{O}_X$-Módulo coherente.
Entonces, para todo $x\in X$, se tiene
\begin{equation}
\label{IV.10.8.3.1}
\tag{10.8.3.1}
\operatorname{prof}^*_x(\sh{F})=\dim(X)-\dim.\operatorname{proj}(\sh{F}_x).
\end{equation}
\end{lemma}

\begin{proof}
En efecto, como $X$ es biequidimensional, se tiene \sref[0]{0.14.3.5.1}
\[
  \dim(\overline{\{x\}})=\dim(X)-\operatorname{codim}(\overline{\{x\}},X)=\dim(X)-\dim(\sh{O}_{X,x})
\]
en virtud de \sref{IV.5.1.2}.
Por otra parte, como $X$ es regular, se tiene por \sref[0]{0.17.3.4}
\[
  \operatorname{prof}(\sh{F}_x)=\dim(\sh{O}_{X,x})-\dim.\operatorname{proj}(\sh{F}_x)
\]
de donde el lema.
\end{proof}

\begin{corollary}[10.8.4]
\label{IV.10.8.4}
Bajo las hipótesis de \sref{IV.10.8.3}, la función $x\mapsto\operatorname{prof}^*_x(\sh{F})$ es semicontinua inferiormente.
\end{corollary}

\begin{proof}
Esto se deduce de \sref{IV.10.8.3.1}, pues $x\mapsto\dim.\operatorname{proj}(\sh{F}_x)$ es semicontinua superiormente \sref{IV.6.11.1}.
\end{proof}

\begin{proposition}[10.8.5]
\label{IV.10.8.5}
Sea $S$ un preesquema localmente noetheriano, $X$ un preesquema localmente de tipo finito sobre $S$ y $\sh{F}$ un $\sh{O}_X$-módulo coherente.
Supongamos que $S$ satisface las condiciones: $1^\circ$~$S$ es un preesquema de Jacobson; $2^\circ$~$S$ es regular; $3^\circ$~las componentes irreducibles de $S$ son equicodimensionales.
Entonces la función $x\mapsto\operatorname{prof}^*_x(\sh{F})$ es semicontinua inferiormente; dicho de otro modo, para todo entero $n$, el conjunto $U_n$ de los $x\in X$ tales que $\operatorname{prof}^*_x(\sh{F})\geq n$ es abierto.
\end{proposition}

\begin{proof}
Como los anillos locales $\sh{O}_{S,s}$ de $S$ son regulares, son universalmente catenarios \sref{IV.5.6.4}; dicho de otro modo, $S$ satisface las condiciones $1^\circ$, $2^\circ$ y $3^\circ$ de \sref{IV.10.6.1}, y por tanto lo mismo ocurre con $X$ \sref{IV.10.6.1},~(i).
Como la noción de profundidad rectificada es entonces de carácter local \sref{IV.10.8.2},~(iii), podemos limitarnos al caso en que $S=\operatorname{Spec}(A)$ y $X=\operatorname{Spec}(B)$ son afines, siendo $A$ un anillo regular y $B$ una $A$-álgebra de tipo finito, y por tanto un cociente de un anillo de polinomios $C=A[T_1,\ldots,T_n]$; este último anillo es regular \sref[0]{0.17.3.7}.
Podemos pues suponer que $X$ es un subpreesquema cerrado de un preesquema regular $Y$ que satisface también las condiciones $1^\circ$, $2^\circ$ y $3^\circ$ de \sref{IV.10.6.1}; teniendo en cuenta la observación \sref{IV.10.8.2},~(ii), quedamos así reducidos al caso en que $X$ es además regular y noetheriano.
Pero como los anillos locales de $X$ son entonces íntegros, las componentes irreducibles de $X$ son
\oldpage[IV-3]{112}
abiertas \sref[I]{I.6.1.10}, y por tanto también podemos suponer que $X$ es irreducible.
Entonces, como los anillos locales de $X$ son catenarios \sref[0]{0.16.5.12}, la hipótesis $3^\circ$ de \sref{IV.10.6.1} implica que $X$ es biequidimensional (\sref{IV.5.1.5} y \sref[0]{0.14.3.3}); basta pues aplicar \sref{IV.10.8.4}.
\end{proof}

Obsérvese que, si $S$ es el espectro de un cuerpo o de $\mathbf{Z}$, satisface las condiciones de \sref{IV.10.8.5}.

\begin{corollary}[10.8.6]
\label{IV.10.8.6}
Bajo las hipótesis de \sref{IV.10.8.5}, para todo $x\in X$, el número $\operatorname{prof}^*_x(\sh{F})$ es el único entero $n$ que tiene la propiedad siguiente: existe un entorno abierto $U$ de $x$ tal que, para todo punto $x'\in U\cap\overline{\{x\}}$ cerrado en $U$, se tiene $\operatorname{prof}(\sh{F}_{x'})=n$.
En particular, para que $\operatorname{prof}^*_x(\sh{F})\geq m$ (resp.\ $\operatorname{prof}^*_x(\sh{F})\leq m$), es necesario y suficiente que exista un entorno abierto $V$ de $x$ en $X$ tal que, para todo punto $x'\in V\cap\overline{\{x\}}$ cerrado en $V$, se tenga $\operatorname{prof}(\sh{F}_{x'})\geq m$ (resp.\ $\operatorname{prof}(\sh{F}_{x'})\leq m$).
\end{corollary}

\begin{proof}
En efecto, podemos limitarnos al caso en que $x$ no es un punto cerrado; si $\operatorname{prof}^*_x(\sh{F})=n$, el conjunto de los $y\in X$ tales que $\operatorname{prof}^*_y(\sh{F})\leq n$ es cerrado en virtud de \sref{IV.10.8.5}, y por tanto contiene a $\overline{\{x\}}$; en virtud de la semicontinuidad inferior de $\operatorname{prof}^*_y(\sh{F})$, existe un entorno abierto $U$ de $x$ tal que $\operatorname{prof}^*_y(\sh{F})=n$ para todo $y\in U\cap\overline{\{x\}}$, y por tanto $\operatorname{prof}(\sh{F}_{x'})=n$ si $x'\in U\cap\overline{\{x\}}$ es cerrado (puesto que la noción de profundidad rectificada es local).
\end{proof}

Para los preesquemas que satisfacen las hipótesis de \sref{IV.10.8.5}, la noción de profundidad rectificada puede por tanto definirse mediante los valores de la profundidad en los \emph{puntos cerrados} de $X$ (estos forman un conjunto muy denso en todo cerrado de $X$).

\begin{proposition}[10.8.7]
\label{IV.10.8.7}
Sea $S$ un preesquema localmente noetheriano que satisfaga las condiciones $1^\circ$, $2^\circ$ y $3^\circ$ de \sref{IV.10.6.1}.
Sea $X$ un preesquema localmente de tipo finito sobre $S$, $\sh{F}$ un $\sh{O}_X$-Módulo coherente y $Y=\operatorname{Supp}(\sh{F})$.
Entonces, para todo $x\in Y$, se tiene
\begin{equation}
\label{IV.10.8.7.1}
\tag{10.8.7.1}
\operatorname{prof}^*_x(\sh{F})=\dim_x(Y)-\operatorname{coprof}(\sh{F}_x).
\end{equation}
\end{proposition}

\begin{proof}
En efecto, por definición, $\operatorname{coprof}(\sh{F}_x)=\dim(\sh{F}_x)-\operatorname{prof}(\sh{F}_x)$, y se deduce de \sref{IV.10.6.4} que $\dim_x(Y)=\dim(\overline{\{x\}})+\dim(\sh{F}_x)$; de donde \sref{IV.10.8.7.1} por definición de $\operatorname{prof}^*_x(\sh{F})$.
\end{proof}

\begin{corollary}[10.8.8]
\label{IV.10.8.8}
Supongamos que las hipótesis sobre $f$ y $\sh{F}$ son las de \sref{IV.9.9.1}. Entonces la función $x\mapsto\operatorname{prof}^*_x(\sh{F}_{f(x)})$ es construible.
\end{corollary}

\begin{proof}
Para todo $s\in S$, la fibra $X_s$, al ser localmente de tipo finito sobre $k(s)$, es un preesquema de Jacobson \sref{IV.10.4.7}. Como $\operatorname{Spec}(k(s))$ satisface las condiciones de \sref{IV.10.6.1}, se tiene
\[
  \operatorname{prof}^*_x(\sh{F}_{f(x)})
  =\dim_x(\operatorname{Supp}(\sh{F}_{f(x)}))
   -\operatorname{coprof}((\sh{F}_{f(x)})_x)
\]
por \sref{IV.10.8.7}.
Ahora bien, si $Z=\operatorname{Supp}(\sh{F})$, se tiene $Z_{f(x)}=\operatorname{Supp}(\sh{F}_{f(x)})$ \sref[I]{I.9.1.13} y $Z$ es localmente construible \sref{IV.8.9.1}; por tanto, las funciones $x\mapsto\dim_x(\operatorname{Supp}(\sh{F}_{f(x)}))$ y $x\mapsto\operatorname{coprof}((\sh{F}_{f(x)})_x)$ son localmente construibles (\sref{IV.9.9.1} y \sref{IV.9.9.3}), lo que prueba la proposición.
\end{proof}

\subsection{Espectros maximales y ultrapreesquemas}
\label{subsection:IV.10.9}

Los resultados de esta sección no se utilizarán en lo que sigue.

\begin{env}[10.9.1]
\label{IV.10.9.1}
Sea $X$ un preesquema de Jacobson, y sea $S(X)$ el \emph{espacio anillado} cuyo espacio subyacente es el subespacio de los puntos cerrados de $X$, y cuyo haz de anillos es el haz inducido sobre este subespacio por $\sh{O}_X$, es decir, el haz $\psi^*(\sh{O}_X)$, donde $\psi:S(X)\to X$ denota la inyección canónica.
Como $\psi$ es un
\oldpage[IV-3]{113}
cuasihomeomorfismo, hemos visto \sref{IV.10.2.8},~(ii) que, si $\theta:\sh{O}_X\to\psi_*(\sh{O}_{S(X)})$ es el homomorfismo de haces de anillos tal que $\theta^\sharp:\psi^*(\sh{O}_X)\to\sh{O}_{S(X)}$ es la identidad, entonces $j_X=(\psi,\theta)$ es un \emph{cuasiisomorfismo} de espacios anillados, y $\sh{F}\mapsto j_X^*(\sh{F})$ es una equivalencia entre la categoría de $\sh{O}_X$-Módulos y la de $\sh{O}_{S(X)}$-Módulos.
Está claro que, bajo esta equivalencia, a los $\sh{O}_X$-módulos localmente libres (resp.\ coherentes) corresponden los $\sh{O}_{S(X)}$-módulos localmente libres (resp.\ coherentes). Además, si $\sh{L}$ es un $\sh{O}_X$-módulo localmente libre y $U$ es un abierto de $X$ tal que $\sh{L}|_U\simeq(\sh{O}_X|_U)^n$, entonces
\[
  j_X^*(\sh{L})|_{U\cap S(X)}
  \simeq \bigl(\sh{O}_{S(X)}|_{U\cap S(X)}\bigr)^n.
\]
\end{env}

\begin{env}[10.9.2]
\label{IV.10.9.2}
Sean $X$, $Y$ dos preesquemas de Jacobson y $f=(\rho,\lambda):X\to Y$ un morfismo \emph{localmente de tipo finito}.
Se ha visto \sref{IV.10.4.7} que $\rho(S(X))\subset S(Y)$ y, por restricción de $\rho$ a $S(X)$, se define por tanto una aplicación continua $S(\rho):S(X)\to S(Y)$.
Para todo abierto $V$ de $Y$, se define por composición un homomorfismo de anillos
\[
  \Gamma(V\cap S(Y),\sh{O}_{S(Y)})\to\Gamma(V,\sh{O}_Y)\xrightarrow{\lambda_V}\Gamma(f^{-1}(V),\sh{O}_X)\to\Gamma(f^{-1}(V)\cap S(X),\sh{O}_{S(X)})
\]
donde los dos isomorfismos extremos se han definido en \sref{IV.10.9.1}; está claro que esto define un homomorfismo de haces de anillos $S(\lambda):\sh{O}_{S(Y)}\to\rho_*(\sh{O}_{S(X)})$ (recordando que los abiertos de $S(X)$ (resp.\ $S(Y)$) corresponden biyectivamente a los de $X$ (resp.\ $Y$) \sref{IV.10.2.1}); se obtiene así un morfismo de espacios anillados $S(f)=(S(\rho),S(\lambda)):(S(X),\sh{O}_{S(X)})\to(S(Y),\sh{O}_{S(Y)})$ tal que el diagrama
\[
  \xymatrix{
    S(X)\ar[r]^{S(f)}\ar[d]_{j_X} & S(Y)\ar[d]^{j_Y}\\
    X\ar[r]^{f} & Y
  }
\]
es conmutativo; además, si $Z$ es un tercer preesquema de Jacobson y $g:Y\to Z$ un morfismo localmente de tipo finito, $S(g\circ f)=S(g)\circ S(f)$.
Se ha definido así un \emph{funtor covariante} $S:\mathbf{C}\to\mathbf{C}'$, donde $\mathbf{C}'$ es la categoría de los \emph{espacios anillados en anillos locales}, y $\mathbf{C}$ la categoría cuyos objetos son los \emph{preesquemas de Jacobson} y cuyos morfismos son los morfismos localmente de tipo finito entre preesquemas de Jacobson.
\end{env}

\begin{env}[10.9.3]
\label{IV.10.9.3}
Determinemos ahora la subcategoría $\mathbf{C}''$ de $\mathbf{C}'$ formada por los espacios anillados isomorfos a $S(X)$ y cuyos morfismos provienen de los $S(f)$.
Supongamos primero que $X=\operatorname{Spec}(A)$, donde $A$ es un anillo de Jacobson; entonces $S(X)$ es el conjunto de los ideales \emph{maximales} de $A$, provisto: $1^\circ$~de la topología inducida por la de $X$, de modo que una base de esta topología está formada por los $D^m(h)=D(h)\cap S(X)$, es decir, el conjunto de los ideales maximales $\mathfrak{m}$ de $A$ tales que $h\notin\mathfrak{m}$; $2^\circ$~del haz de anillos $\sh{O}_{S(X)}$ tal que $\Gamma(D^m(h),\sh{O}_{S(X)})=A_h$.
Diremos que este espacio anillado es el \emph{espectro maximal} del anillo de Jacobson $A$ y lo denotaremos $\operatorname{Spm}(A)$.

Obsérvese que, si $j:D(h)\to X$ es la inyección canónica, el espacio anillado inducido sobre $D^m(h)$ por $S(X)$ y $S(D(h))$ son canónicamente isomorfos, y que la inyección canónica $D^m(h)\to S(X)$ de espacios anillados es igual a $S(j)$.

Sea $B$ un segundo anillo de Jacobson, $Y=\operatorname{Spec}(B)$, $\varphi:B\to A$ un homomorfismo de anillos que hace de $A$ una $B$-álgebra de \emph{tipo finito}, $f=({}^a\varphi,\widetilde{\varphi}):X\to Y$ el morfismo correspondiente de preesquemas, y
\[
  S(f):\operatorname{Spm}(A)\to\operatorname{Spm}(B)
\]
el morfismo de espacios anillados correspondiente a $f$.
Está claro que $S(f)=(\psi,\theta)$ es un morfismo de espacios anillados en anillos locales, es decir, $(\mathrm{Err}_{\mathrm{II}},1.8.2)$, para todo $x\in\operatorname{Spm}(A)$, $\theta^\sharp_x$ es un homomorfismo \emph{local}.
Recíprocamente:
\end{env}

\begin{proposition}[10.9.4]
\label{IV.10.9.4}
Sean $A$, $B$ dos anillos de Jacobson.
Si $u=(\psi,\theta):\operatorname{Spm}(A)\to\operatorname{Spm}(B)$ es un morfismo de espacios anillados en anillos locales tal que $\Gamma(\theta):B\to A$ hace de $A$ una $B$-álgebra de tipo finito, existe un único morfismo de preesquemas $f:\operatorname{Spec}(A)\to\operatorname{Spec}(B)$ tal que $u=S(f)$.
\end{proposition}

\begin{proof}
La unicidad de $f$ es evidente, pues si $f=({}^a\varphi,\widetilde{\varphi})$, se debe tener $\varphi=\Gamma(\theta)$; queda por mostrar que $S(f)$ está definido y que $u=S(f)$.
Ahora bien, la primera afirmación se deduce de que se supone que $\varphi$ hace de $A$ una $B$-álgebra de tipo finito, y por tanto $f(\operatorname{Spm}(A))\subset\operatorname{Spm}(B)$; el hecho de que $\theta^\sharp_x$ sea un homomorfismo local para todo $x\in\operatorname{Spm}(A)$ permite repetir el razonamiento de \sref[I]{I.1.7.3}, restringiéndonos a los puntos $x$ de $\operatorname{Spm}(A)$, para mostrar así sucesivamente que $\psi(x)={}^a\varphi(x)$ para todo $x\in\operatorname{Spm}(A)$, y después que $\theta^\sharp_x=\varphi_x:B_{\psi(x)}\to A_x$, lo que termina de probar que $u=S(f)$.
\end{proof}

\oldpage[IV-3]{114}

\begin{env}[10.9.5]
\label{IV.10.9.5}
Consideremos ahora un espacio anillado $(X,\sh{O}_X)$; diremos que un abierto $U$ de $X$ es \emph{ultraafín} si el espacio anillado inducido $(U,\sh{O}_X|U)$ es isomorfo a un espectro maximal $\operatorname{Spm}(A)$, donde $A$ es un anillo de Jacobson.
Diremos que $X$ es un \emph{ultrapreesquema} si todo punto de $X$ admite un entorno abierto ultraafín.
Se muestra, como en \sref[I]{I.2.1.3} y \sref[I]{I.2.1.4}, que los abiertos ultraafines forman una \emph{base} de la topología de $X$ y que $X$ es un espacio de Kolmogórov.
Si $Y$ es un segundo ultrapreesquema, diremos que un morfismo de espacios anillados $f:X\to Y$ es un \emph{morfismo de ultrapreesquemas} si satisface las condiciones siguientes: $1^\circ$~$f$ es un morfismo de espacios anillados en anillos locales; $2^\circ$~para todo $x\in X$, existe un entorno abierto ultraafín $V$ de $f(x)$ en $Y$ y un abierto ultraafín $U$ de $x$ en $X$ tales que $f(U)\subset V$ y que el homomorfismo $\Gamma(V,\sh{O}_Y)\to\Gamma(U,\sh{O}_X)$ correspondiente a $f$ hace de $\Gamma(U,\sh{O}_X)$ una $\Gamma(V,\sh{O}_Y)$-\emph{álgebra de tipo finito}.

Es inmediato que la compuesta de dos morfismos tales está bien definida y es un tercer morfismo tal por la observación final de \sref{IV.10.9.3}.
Está claro que la categoría $C''_0$ así definida es una subcategoría de $C'$ que contiene a $C''$; nos proponemos ahora mostrar que $C''=C''_0$, dicho de otro modo:
\end{env}

\begin{proposition}[10.9.6]
\label{IV.10.9.6}
El funtor $X\mapsto S(X)$ de $\mathbf{C}$ en $\mathbf{C}''_0$ es una equivalencia de categorías.
\end{proposition}

\begin{proof}
$1^\circ$ Mostremos primero que el funtor $X\mapsto S(X)$ es \emph{plenamente fiel}, dicho de otro modo, que para $X$, $Y$ en $\mathbf{C}$, la aplicación canónica
\[
  \operatorname{Hom}_{\mathbf{C}}(X,Y)\to\operatorname{Hom}_{\mathbf{C}''_0}(S(X),S(Y))
\]
es \emph{biyectiva}.
En primer lugar es \emph{inyectiva}: sean $f$, $g$ dos morfismos localmente de tipo finito de $X$ en $Y$, y supongamos que $S(f)=S(g)$; se tiene $f^{-1}(V)\cap S(X)=g^{-1}(V)\cap S(X)$ para todo abierto $V$ de $Y$, luego \sref{IV.10.3.1} y \sref{IV.10.2.7} dan $f^{-1}(V)=g^{-1}(V)$; basta por tanto probar que, para todo abierto afín $V$ de $Y$, $f$ y $g$ coinciden en $f^{-1}(V)=g^{-1}(V)$; dicho de otro modo, nos reducimos al caso en que $Y=\operatorname{Spec}(B)$ es el espectro de un anillo de Jacobson $B$.
Basta \sref[I]{I.2.2.4} mostrar que los homomorfismos de anillos $B\to\Gamma(X,\sh{O}_X)$ correspondientes a $f$ y $g$ son iguales.
Ahora bien, para todo $s\in B$, las imágenes de $s$ por estos homomorfismos son dos secciones de $\sh{O}_X$ sobre $X$ que, por hipótesis, inducen la \emph{misma} sección sobre $S(X)$; sabemos \sref{IV.10.2.8} que estas secciones son idénticas, de donde nuestra afirmación.
En segundo lugar, probemos que todo morfismo $h:S(X)\to S(Y)$ (para la categoría $\mathbf{C}''_0$) es de la forma $S(f)$, donde $f:X\to Y$ es un morfismo localmente de tipo finito.
Por hipótesis, existe un recubrimiento abierto ultraafín $(U'_\lambda)$ (resp.\ $(V'_\mu)$) de $S(X)$ (resp.\ $S(Y)$) tal que, para todo $\lambda$, $h(U'_\lambda)$ está contenido en uno de los $V'_\mu$ y el homomorfismo de anillos correspondiente a la restricción $h_\lambda:U'_\lambda\to V'_\mu$ hace del segundo anillo un álgebra de tipo finito sobre el primero.
Podemos suponer que $U'_\lambda=U_\lambda\cap S(X)$ y $V'_\mu=V_\mu\cap S(Y)$, donde $U_\lambda$ y $V_\mu$ son abiertos afines de $X$ y $Y$, respectivamente, determinados de manera única. Si, para todo $\lambda$, $h_\lambda=S(f_\lambda)$ para un morfismo de tipo finito $f_\lambda:U_\lambda\to V_\mu$, entonces la primera parte del razonamiento, aplicada a las restricciones de $f_\alpha$ y $f_\beta$ a $U_\alpha\cap U_\beta$, muestra que los $f_\lambda$ son restricciones de un único morfismo $f:X\to Y$, y claramente $h=S(f)$. Quedamos así reducidos al caso en que $X$ e $Y$ son afines, y la conclusión se deduce de \sref{IV.10.9.4}.

$2^\circ$ Queda por probar que todo ultrapreesquema $X'$ es de la forma $S(X)$ para un preesquema de Jacobson $X$ (que será necesariamente único salvo isomorfismo, en virtud de $1^\circ$).
Existen abiertos ultraafines $(U'_\alpha)$ de $X'$, cada uno de los cuales es de la forma $S(U_\alpha)$, siendo $U_\alpha$ el espectro de un anillo de Jacobson.
Para cada par de índices $\alpha$, $\beta$, consideremos el único abierto $U_{\alpha\beta}$ de $U_\alpha$ cuya traza sobre $U'_\alpha$ es $U'_\alpha\cap U'_\beta$; en virtud de $1^\circ$, el automorfismo identidad de $U'_\alpha\cap U'_\beta$ es de la forma $S(\theta_{\alpha\beta})$, donde $\theta_{\alpha\beta}:U_{\alpha\beta}\to U_{\beta\alpha}$ es un isomorfismo de preesquemas.
Se verifica fácilmente (en virtud de $1^\circ$) que la familia $(\theta_{\alpha\beta})$ satisface la condición de pegado \sref[0]{0.4.1.7}, y que esta familia define un preesquema $X$; está claro entonces que se tiene $X'=S(X)$, lo que termina la demostración.
\end{proof}

\subsection{Espacios algebraicos de Serre}
\label{subsection:IV.10.10}

\begin{env}[10.10.1]
\label{IV.10.10.1}
El lenguaje introducido por Serre en (FAC) resulta a veces cómodo, especialmente en las cuestiones en que el interés principal recae sobre los puntos racionales sobre el cuerpo base $k$ (algebraicamente cerrado por hipótesis) de las ``variedades algebraicas'' sobre $k$ que se consideran.
Esbozaremos aquí este lenguaje enlazándolo con las consideraciones precedentes, para permitir al lector traducir los enunciados de Serre al lenguaje de los esquemas.
Por lo demás, también es posible desarrollar el lenguaje de Serre para preesquemas sobre un cuerpo no algebraicamente cerrado (e incluso sobre un anillo artiniano); pero esto introduce complicaciones técnicas considerables y, además, sobre un cuerpo base arbitrario desaparecen las ventajas (ante todo psicológicas) del punto de vista de Serre; por ello nos limitaremos al marco fijado por Serre.
La presente sección, al igual que la precedente, no se utilizará en lo que sigue de este Tratado, y nos contentaremos por tanto con indicaciones breves.
\end{env}

\oldpage[IV-3]{115}

\begin{env}[10.10.2]
\label{IV.10.10.2}
Dado un ultrapreesquema $R$, se puede definir naturalmente (como en toda categoría) la noción de $R$-\emph{ultrapreesquema}.
Consideremos en particular un cuerpo algebraicamente cerrado $k$; $\operatorname{Spm}(k)$ es entonces idéntico a $\operatorname{Spec}(k)$; diremos que un $\operatorname{Spm}(k)$-ultrapreesquema $X'$ es un \emph{espacio prealgebraico sobre $k$}: es un espacio $k$-anillado en anillos locales en el que todo punto admite un entorno abierto isomorfo al espectro maximal de una $k$-álgebra de tipo finito; por tanto, esto equivale a decir (por \sref{IV.10.9.6}) que $X'=S(X)$, donde $X$ es un $k$-preesquema localmente de tipo finito.
Si $X$, $Y$, $Z$ son tres $k$-preesquemas localmente de tipo finito, lo mismo ocurre con $X\times_ZY$ \sref[I]{subsection:I.3.4}; por tanto, en virtud de \sref{IV.10.9.6}, existen productos en la categoría de los espacios prealgebraicos sobre $k$, tanto el producto ``sobre $k$'' $X'\times_kY'$ como el ``producto fibrado'' $X'\times_{Z'}Y'$, donde $X'$, $Y'$, $Z'$ son tres espacios prealgebraicos sobre $k$.
Se puede definir, pues, el morfismo diagonal $\Delta_{X'}:X'\to X'\times_k X'$ (que no es otro que $S(\Delta_X)$); se dice que $X'$ es un \emph{espacio algebraico sobre $k$} si $X$ es un esquema, y esto equivale a decir que la imagen de $\Delta_{X'}$ es un subconjunto \emph{cerrado} de $X'\times_k X'$.
\end{env}

\begin{env}[10.10.3]
\label{IV.10.10.3}
Las simplificaciones que provienen de la hipótesis de que $k$ es algebraicamente cerrado consisten ante todo en lo siguiente: para todo $k$-preesquema $X$ localmente de tipo finito, existe una correspondencia biyectiva entre los \emph{puntos cerrados de $X$}, los \emph{puntos de $X$ con valores en $k$} \sref[I]{I.3.4.4} y los \emph{puntos racionales de $X$ sobre $k$} \sref[I]{I.3.4.5}, en virtud de \sref[I]{I.6.4.2}.
Esto muestra en particular \sref[I]{I.3.4.3.1} que, para dos espacios $k$-prealgebraicos $X'$, $Y'$, el conjunto subyacente al producto $X'\times_kY'$ es idéntico al \emph{conjunto producto} $X'\times Y'$ de los conjuntos subyacentes (pero, desde luego, la topología del espacio subyacente a $X'\times_kY'$ no es la \emph{topología producto} de las topologías de los espacios subyacentes a $X'$ e $Y'$ y, en general, es estrictamente más fina que esta última).

Por otra parte, los anillos locales $\sh{O}_x$ en los puntos de un espacio prealgebraico $X'$ sobre $k$ son $k$-álgebras cuyo cuerpo residual es \emph{isomorfo a $k$}; si $A$ y $B$ son dos $k$-álgebras de esta índole, un $k$-homomorfismo $\varphi:A\to B$ es necesariamente \emph{local}: en efecto, si existiera un elemento $x$ del ideal maximal de $A$ tal que $\varphi(x)$ fuese invertible, existiría un $\lambda\in k$ no nulo tal que $\varphi(x-\lambda.1)$ perteneciese al ideal maximal de $B$, lo cual es absurdo puesto que $x-\lambda.1$ es invertible.
De ello concluimos inmediatamente que, si $X'=S(X)$, $Y'=S(Y)$ son dos espacios prealgebraicos sobre $k$, todo morfismo $X'\to Y'$ de \emph{espacios $k$-anillados} es también un morfismo de espacios $k$-anillados \emph{en anillos locales} (\sref[I]{I.1.8.2}; cf.\ $\mathrm{Err}_{\mathrm{II}}$). Además, con la notación precedente, si $A$ y $B$ son $k$-álgebras de tipo finito, el homomorfismo correspondiente hace de $B$ una $A$-álgebra de tipo finito; por tanto, todo morfismo $X'\to Y'$ de espacios $k$-anillados es, por \sref{IV.10.9.6}, de la forma $S(f)$, donde $f:X\to Y$ es un morfismo de $k$-preesquemas.

Por último, para todo abierto $U$ de $X'$, toda sección $s\in\Gamma(U,\sh{O}_{X'})$ y todo $x\in U$, $s(x)$ \sref[0]{0.5.5.1} se identifica con un elemento de $k$; así, $s$ determina una aplicación $\overline{s}:x\mapsto s(x)$ de $U$ en $k$, es decir, una sección sobre $U$ del haz $\sh{A}(X')$ de gérmenes de aplicaciones de $X'$ en $k$. Como $h_U:s\mapsto\overline{s}$ es evidentemente un homomorfismo de anillos $\Gamma(U,\sh{O}_{X'})\to\Gamma(U,\sh{A}(X'))$ y conmuta con la restricción a un abierto $V\subset U$, los $h_U$ definen un homomorfismo de haces de anillos $h:\sh{O}_{X'}\to\sh{A}(X')$.
Si se toma $U$ ultraafín, $\operatorname{Spm}(A)$, donde $A$ es un anillo de Jacobson, decir que $\overline{s}=0$ significa que, para todo ideal maximal $\mathfrak{m}$ de $A$, $s$ pertenece a $\mathfrak{m}$, o equivalentemente que $s$ pertenece al \emph{radical} de $A$; pero, como $A$ es un anillo de Jacobson, su radical es su nilradical. Por consiguiente, $h_U$ es inyectivo si y solo si $A$ es reducido.
\end{env}

\begin{env}[10.10.4]
\label{IV.10.10.4}
Se dice que el espacio $k$-prealgebraico $X'=S(X)$ es \emph{reducido} si $X$ es reducido. Como el conjunto de puntos $x\in X$ en los que $X$ es reducido es abierto \sref[0]{0.5.2.2}, su complementario, si no es vacío, contiene al menos un punto cerrado \sref{IV.5.1.11}; por tanto, equivale a decir que $X'$ es reducido o que cada uno de sus anillos locales $\sh{O}_x$ ($x\in X'$) es reducido. Acabamos de ver en \sref{IV.10.10.3} que, para que el homomorfismo $h:\sh{O}_{X'}\to\sh{A}(X')$ sea \emph{inyectivo}, es necesario y suficiente que $X'$ sea \emph{reducido}.
En (FAC), Serre se limita de hecho a los espacios prealgebraicos \emph{reducidos}, lo que le permite definir $\sh{O}_{X'}$ como un \emph{subhaz} de $\sh{A}(X')$.
Obsérvese que, si $X'$ e $Y'$ son dos espacios $k$-prealgebraicos reducidos, lo mismo ocurre con $X'\times_k Y'$: en efecto, esto equivale a decir que, si $A$ y $B$ son dos $k$-álgebras reducidas de tipo finito, $A\otimes_k B$ es también reducido; pero acabamos de ver que los radicales de $A$ y $B$ son $0$, y como $k$ es algebraicamente cerrado, $A$ y $B$ son álgebras ``separables'' sobre $k$ en el sentido de Bourbaki (\emph{Alg.}, cap.~VIII, \S\,7, n$^\circ$\,5, prop.~5); por tanto, $A\otimes_k B$ carece de radical (\emph{loc.\ cit.}, n$^\circ$\,6, cor.~3 del th.~3), y como es un anillo de Jacobson, es reducido.
Sin embargo, si $Z'$ es un tercer espacio prealgebraico reducido sobre $k$, el ``producto fibrado'' $X'\times_{Z'}Y'$ de dos espacios prealgebraicos reducidos sobre $Z'$ no es reducido en general, lo que implica que la categoría de estos espacios resulta insuficiente en muchas cuestiones (especialmente en la teoría de grupos algebraicos).
Pero, como hemos visto anteriormente, se puede conservar el lenguaje de Serre sin limitarse, como hace este último (quien además solo considera espacios \emph{prealgebraicos} cuasicompactos), al caso de los espacios prealgebraicos reducidos.
\end{env}

\oldpage[IV-3]{116}

\begin{env}[10.10.5]
\label{IV.10.10.5}
Por último, también se pueden considerar ultrapreesquemas sobre un cuerpo arbitrario $k$
conservando un lenguaje próximo al de Serre e introduciendo, como en Weil, una
extensión algebraicamente cerrada fija $K$ de $k$ (elegida suficientemente grande---por ejemplo,
de grado de trascendencia infinito sobre $k$---para disponer de suficientes ``puntos genéricos''
en el sentido de Weil).
A todo preesquema $X$ localmente de tipo finito sobre $k$ se asocia entonces el conjunto
\[
 S_K(X)=S(X\otimes_k K)
\]
de puntos de $X$ con valores en $K$.
Existe una aplicación canónica $j:S_K(X)\to X$, que es un cuasihomeomorfismo cuando
se dota a $S_K(X)$ de la topología imagen inversa de la de $X$ por $j$; al dotar a $S_K(X)$
del haz de $k$-álgebras $j^*(\sh{O}_X)$, se obtiene una subcategoría de la categoría
de espacios $k$-anillados en anillos locales, que se podría llamar categoría de los
\emph{espacios $(k,K)$-prealgebraicos}.
Todavía se pueden definir productos en ella, y los resultados de \sref{IV.10.10.3} y
\sref{IV.10.10.4} se pueden generalizar (sustituyendo aquí $\sh{A}(X')$ por el haz
$\sh{A}_K(X')$ de gérmenes de aplicaciones de $X'$ en $K$).
Sin embargo, este punto de vista atribuye un papel artificial a un cuerpo de extensión
de $k$ elegido arbitrariamente, y lo mencionamos únicamente para descartarlo.
\end{env}
